% =============================================================================
% manuscript_zh.tex — 论文中文逐句对照版（人工审稿用）
% 对应: ai_submission/manuscript.tex (commit 060e2cf)
% 编译: xelatex manuscript_zh.tex （xelatex 两遍）
% 说明: 正文叙述逐句译为中文，数学公式/符号/表格数据/参考文献原文保留，
%       定理、定义、命题、引理、注记、证明环境改为中文标题，标签与原文一致，
%       便于与英文版逐句对照。
% =============================================================================
% 注意: 本机 TeX Live 2026basic 未安装 ctex/xeCJK，故用纯 fontspec + XeTeX
%       内置中文断行（\XeTeXlinebreaklocale）实现，主字体用 macOS 系统字体。
% =============================================================================
\documentclass[11pt,a4paper]{article}
\usepackage{amsmath,amssymb,amsfonts}
\usepackage{booktabs}
\usepackage{multirow}
\usepackage{graphicx}
\usepackage{url}
\usepackage{geometry}
\geometry{margin=2.4cm}
\usepackage{amsthm}
\usepackage{fontspec}
\setmainfont{Songti SC}
\XeTeXlinebreaklocale "zh"
\XeTeXlinebreakskip = 0pt plus 1pt

\newtheorem{theorem}{定理}[section]
\newtheorem{lemma}[theorem]{引理}
\newtheorem{proposition}[theorem]{命题}
\newtheorem{corollary}[theorem]{推论}
\newtheorem{definition}[theorem]{定义}
\newtheorem{remark}[theorem]{注记}
\renewcommand{\proofname}{证明}

\newcommand{\coloneqq}{\triangleq}

\title{类级机器学习遗忘的谱认证（逐句对照中文版）}
\author{Yan Zhang \and Chenggang Lu\\[4pt] {\small 张炎、卢成刚（浙江工业大学理学院）}}
\date{2026-08-13\quad 对照版本（commit 060e2cf）}

\begin{document}

\maketitle

% =============================================================================
% 摘要
% =============================================================================
\begin{abstract}
\noindent
\textbf{摘要（逐句对照）：}
机器学习遗忘旨在无需完全重训练的情况下，从已训练模型中移除特定训练数据的影响。验证这种移除是否有效仍是一个开放且具有实际重要性的问题。成员推理攻击（MIA）将训练/测试分布位移与遗忘集泄漏混为一谈，而差分隐私界要求侵入式训练，因此二者都无法证明输出级不可区分性。我们引入一个基于 $2\times C$ \emph{经验输出通道矩阵}的谱证书——两行分别是遗忘集与保留集上的平均 softmax 预测——并定义 \textbf{残差不可逆指数（Residual Irreversibility Index, RII）} $\rho=1-\sigma_1^2/(\sigma_1^2+\sigma_2^2)$，当两个通道输出分布不可区分时它精确地消失。对于类级遗忘（标准 RII 在此遭受已知/未知类不对称），我们引入 \textbf{多留出投影残差（Multi-Held-Out Projection Residual, MHPR）} $\rho_H$，将遗忘类均值投影到多个留出类的子空间上，并提供有限样本保证与温度缩放。我们证明 $\rho=0$ 等价于输出通道独立，并导出微调泄漏、互信息与统一两轴分解界。在 CIFAR-10、Fashion-MNIST、CIFAR-100 以及 TOFU 基准~\cite{maini2024tofu} 与 MUSE 风格书籍基准上，配合 70 亿参数语言模型（LLaMA-2-7B，LoRA），RII 在大多数设置中对方法的排序与已知遗忘强度一致，并在准确率与 MIA 饱和处仍具信息量。当生成准确率饱和在 100\% 时，准确率本身无法对检查点排序，而 RII 仍将微调排序在未遗忘之下。RII 揭示出梯度上升的一个隐藏失败（其输出签名在准确率与 MIA 改善的同时反而增强）以及一个基于准确率的指标无法发现的微调反弹。在长文本书籍遗忘中，梯度上升同样使遗忘输出分布保持不变，遗忘困惑度将重训练排序为最强擦除，而原始词表通道反映的是书籍的固有差异而非遗忘状态，界定了 RII 具有信息量的情形。表示级探针证实输出级擦除可以与特征级签名共存，这促使采用两层审计。

\medskip
\noindent\textbf{关键词：}秩一通道、机器学习遗忘、奇异值分解、残差不可逆指数、多留出投影、输出空间不可逆性
\end{abstract}

% =============================================================================
% §1 引言
% =============================================================================
\section{引言}
\label{sec:intro}

被遗忘权（GDPR 第 17 条）要求机器学习模型应请求消除特定训练数据的影响~\cite{bourtoule2021machine,nguyen2025survey}。一个核心挑战是\emph{如何验证遗忘已经成功}。现有方法要么依赖将训练/测试分布位移与遗忘/保留泄漏混为一谈的成员推理攻击（MIA），要么依赖要求侵入式训练的差分隐私界。

具体而言，我们研究的问题是：\emph{能否用一个轻量、黑盒的证书，在类级情形下（即现有指标最不可靠的情形）验证遗忘集与保留集之间的输出级不可区分性？}我们的答案是模型输出通道上的一个谱条件。

我们引入一个基于\textbf{谱条件}的轻量黑盒证书：输出空间不可逆性归结为 $2\times C$ 经验输出通道矩阵的秩。所谓\textbf{输出分布不可逆性}，我们指遗忘集成员指示变量 $X$ 与模型在参数 $\theta'$ 下的类别输出 $Y$ 条件独立，即 $I(X;Y\mid\theta')=0$。当此成立时，模型在 $D_f$ 与 $D_r$ 上的输出分布在统计上不可区分，产生一个秩一输出通道矩阵且 $\rho=0$。

\textbf{范围。}我们的框架严格在\emph{输出空间}（可观测的预测）内度量不可逆性，这与黑盒 GDPR 合规审计相符。我们不声称参数级白盒安全。我们聚焦于\textbf{类级遗忘}——最严格且最具实践相关性的情形，其中 $D_f$ 由整个语义类的全部样本组成；MHPR（见第~\ref{sec:mhpr}~节）专门处理此处出现的已知/未知类不对称。

框架建立在一个基本观察之上。$2\times C$ \textbf{经验输出通道矩阵} $\mathbf{M}(\theta')$ 的行是 $\boldsymbol{\mu}_f$ 与 $\boldsymbol{\mu}_r$（$D_f$ 与 $D_r$ 上的平均 softmax 向量）。其奇异值 $\sigma_1 \ge \sigma_2$ 编码残差依赖：$\sigma_2 = 0 \iff \mathrm{rank}(\mathbf{M}) = 1 \iff$ 完美移除。

\textbf{贡献：}
\begin{enumerate}
\item 一个输出谱框架（第~\ref{sec:channel}~节）：一个 $2\times C$ 通道矩阵，其秩一条件给出无需调参的残差不可逆指数（RII）。
\item 一个证明（定理~\ref{thm:perfect}）：在\emph{通道}层面，$\rho=0$ 当且仅当 $I(X;Y\mid\theta')=0$，无需指数族或校准假设；逐样本输出分布刻意不在范围之内。
\item 面向类级遗忘的 MHPR（第~\ref{sec:mhpr}~节）：一种多留出投影，在很大程度上解决已知/未知类不对称（将 oracle 的残差降到远低于其标准 RII 水平），并附带温度缩放与有限样本保证。
\item 理论保证（第~\ref{sec:theory}~节）：置信界、$O(\lambda^2)$ 微调泄漏界、谱互信息界，以及将输出泄漏分解为两个正交轴的统一分解。
\item 对 RII 与 MHPR 在统一七方法基准以及 TOFU 与 MUSE 风格书籍遗忘基准的 70 亿参数语言模型上的评估（第~\ref{sec:experiments}~节），揭示准确率与 MIA 都会错过的隐藏失败（梯度上升签名增强与微调反弹）；在 TOFU 上，RII 对生成准确率饱和于 100\% 的已遗忘检查点进行排序，而准确率本身不提供任何排序，且在长文本书籍遗忘中梯度上升同样使遗忘输出分布保持不变。我们还给出 $O(CN)$ 审计协议；逐样本精化（PSSD）在附录~\ref{app:pssd} 中给出。
\end{enumerate}

% =============================================================================
% §2 问题设置与谱框架
% =============================================================================
\section{问题设置与谱框架}
\label{sec:channel}

\subsection{问题设置}

设 $D = \{(x_i, y_i)\}_{i=1}^N$ 为训练集，$\theta_0 = \mathcal{A}(D)$ 为训练所得参数，$D_f \subset D$ 为遗忘集，$D_r = D \setminus D_f$ 为保留集，并令 $N_f = |D_f|$、$N_r = |D_r|$。遗忘算法 $\mathcal{U}$ 产生 $\theta' = \mathcal{U}(\theta_0, D_f, D_r)$。我们通过模型诱导的输出分布来评估 $D_f$ 的残差影响。

我们聚焦于\textbf{类级遗忘}，其中 $D_f$ 由整个语义类的全部样本组成。这一情形暴露出样本级遗忘可能掩盖的谱结构。

\subsection{经验输出通道矩阵}

设 $f_{\theta'}$ 为 logit 函数，$\mathbf{p}(x;\theta') = \mathrm{Softmax}(f_{\theta'}(x)) \in \Delta^{C-1}$。所有量都由 softmax 概率向量而非硬（argmax）预测计算，抽象通道变量 $Y$（定理~\ref{thm:perfect}）是从 $p_c(x;\theta')$ 抽取的类；因此该证书涉及概率输出行为，而非标签一致性。

\begin{definition}[经验输出通道矩阵]
\label{def:M}
\begin{equation}
\mathbf{M}(\theta') =
\begin{bmatrix}
\boldsymbol{\mu}_f^\top \\ \boldsymbol{\mu}_r^\top
\end{bmatrix}
\in \mathbb{R}^{2\times C},
\qquad
\begin{aligned}
\boldsymbol{\mu}_f &= \frac{1}{N_f}\sum_{x\in D_f} \mathbf{p}(x;\theta'),\\
\boldsymbol{\mu}_r &= \frac{1}{N_r}\sum_{x\in D_r} \mathbf{p}(x;\theta').
\end{aligned}
\end{equation}
\end{definition}

对定义~\ref{def:M} 的通道矩阵，其 SVD 为 $\mathbf{M} = \sigma_1 \mathbf{u}_1 \mathbf{v}_1^\top + \sigma_2 \mathbf{u}_2 \mathbf{v}_2^\top$，其中 $\sigma_1 \ge \sigma_2 \ge 0$。奇异值平方是 $2\times 2$ Gram 矩阵 $\mathbf{M}\mathbf{M}^\top$ 的特征值；当 $\|\boldsymbol{\mu}_f\| = \|\boldsymbol{\mu}_r\|$ 时，有 $\|\boldsymbol{\mu}_f - \boldsymbol{\mu}_r\|_2 = \sqrt{2}\,\sigma_2$。

\subsection{残差不可逆指数（RII）}

\begin{definition}[RII]
\label{def:rii}
\begin{equation}
\rho(\theta') \coloneqq 1 - \frac{\sigma_1^2}{\sigma_1^2 + \sigma_2^2}
= \frac{\sigma_2^2}{\sigma_1^2 + \sigma_2^2}
\;\in\; [0,\,0.5].
\label{eq:rho}
\end{equation}
\end{definition}

性质：(i) $\rho = 0 \iff \sigma_2 = 0 \iff \boldsymbol{\mu}_f = \boldsymbol{\mu}_r$（完美遗忘）；(ii) 当 $\|\boldsymbol{\mu}_f\|_2=\|\boldsymbol{\mu}_r\|_2$ 且 $\boldsymbol{\mu}_f \perp \boldsymbol{\mu}_r$ 时 $\rho = 0.5$；(iii) 尺度不变；(iv) 当 $\rho \ll 1$ 时 $\rho \approx \sigma_2^2/\sigma_1^2$。

定义~\ref{def:rii}（式~\eqref{eq:rho}）将 $\rho$ 定义为一种谱残差能量（输出通道的归一化第二奇异值），而非互信息值。其与互信息的联系由定理~\ref{thm:mi_bound} 确立：当 $\rho\ll1$ 时 $I(X;Y\mid\theta')\le O(\rho)$，因此 $\rho$ 在互信息消失时精确消失，并在小泄漏时从上方界定它。

\paragraph{统计依据。}
行 $\boldsymbol{\mu}_f, \boldsymbol{\mu}_r$ 是经验蒙特卡洛平均：当 $N_f\to\infty$ 时 $\hat{\boldsymbol{\mu}}_f\to\boldsymbol{\mu}_f^*$，收敛速度 $O_p(1/\sqrt{N_f})$，因此秩一条件是关于\emph{群体级}统计不可区分性的论断（有限样本保证见命题~\ref{prop:confidence}）。

% =============================================================================
% §3 类级遗忘与 MHPR
% =============================================================================
\section{类级遗忘与 MHPR}
\label{sec:class_baselines}

在类级遗忘中，遗忘行与保留行平均的是结构不同的分布：$\boldsymbol{\mu}_f$ 来自模型训练中从未见过的类。因此即使完美遗忘，两行也可能不同，从而抬高标准 RII。

\subsection{单一留出类（LOCO）失效}

用单个留出未见类的均值替换 $\boldsymbol{\mu}_r$（LOCO）会失败，因为不同未见类产生系统不同的预测轮廓：在 MNIST 上（在 $\{1,2,3,4,6,7,8,9\}$ 上训练，遗忘类~5，留出类~0），$\rho_{\mathrm{LOCO}}{\approx}0.38$ 超过标准 RII ${\approx}0.14$，与预期方向相反。

\subsection{多留出投影残差（MHPR）}
\label{sec:mhpr}

MHPR 使用 $K \ge 2$ 个留出未见类，并将遗忘类均值投影到它们的张成子空间上，从而消除单类偏差。

\begin{definition}[MHPR]
\label{def:mhpr}
设 $\mathcal{H} = \{c_{h1},\dots,c_{hK}\}$ 为 $K$ 个留出类，均值向量为 $\boldsymbol{\mu}_{h_k}$，并令 $\mathbf{H} = [\boldsymbol{\mu}_{h_1},\dots,\boldsymbol{\mu}_{h_K}]^\top$。MHPR 定义为：
\begin{equation}
\rho_H(\theta') \coloneqq \frac{\|\boldsymbol{\mu}_f - \mathbf{H}^+ \mathbf{H} \boldsymbol{\mu}_f\|_2^2}{\|\boldsymbol{\mu}_f\|_2^2} \in [0,1].
\end{equation}
\end{definition}

与谱 RII $\rho\in[0,0.5]$（$2\times C$ 通道的归一化第二奇异值）不同，$\rho_H$ 是遗忘类均值在留出子空间上的\emph{投影残差的相对能量}；归一化对象不同（谱能量比 vs.\ 几何残差范数比）正是两者取值区间不同的原因。一个注意点是按 $\|\boldsymbol{\mu}_f\|_2^2$ 归一化：当模型对遗忘类接近均匀（例如重度遗忘或退化下）时，$\|\boldsymbol{\mu}_f\|_2$ 很小，残差比例会放大估计噪声；此时未归一化的谱 RII 是更稳定的量。

\paragraph{性质。}
对定义~\ref{def:mhpr} 的 MHPR：(1) 单调：当 $K_2 \ge K_1$ 时 $\rho_H(K_2) \le \rho_H(K_1)$（命题~\ref{prop:mhpr_monotone}）。(2) 优于 LOCO：当 $K\ge 2$ 时 $\rho_H(K) \le \rho_H(1)$。(3) 偏差--方差：在 $\mathcal{H}_0$（理想遗忘）下，$\mathbb{E}[\rho_H] = O(CK/N_{\min})$（命题~\ref{prop:mhpr_bias}）；当 $K=O(1)$ 时为 $O(C/N_{\min})$。(4) 有限样本界：$|\hat{\rho}_H - \rho_H^*|$ 以 $O(1/\sqrt{N_{\min}})$ 衰减并带显式常数（命题~\ref{prop:mhpr_confidence}）。证明见附录。

\paragraph{温度缩放。}
当 softmax 输出过度集中（准确率 $\lesssim 90\%$）时，留出类均值向量可能近乎共线，将子空间 $\mathcal{S}_K$ 坍缩为一维流形；我们在一次 Fashion-MNIST 运行中观察到这种坍缩（$\rho_H\approx1$），但在三个种子上它是种子依赖的而非系统性的（第~\ref{sec:benchmark}~节）。

为缓解输出集中，我们应用温度缩放 $\mathbf{p}_T(x) = \mathrm{Softmax}(f_{\theta'}(x)/T)$：当 $T\to\infty$ 时，$\mathbf{p}_T \to \mathbf{1}/C$（均匀），留出子空间变得良态。我们选择条件数满足 $\kappa(\mathbf{H}(T)) < 10$ 的最小 $T\in\{1,2,5,10,50,100\}$。在 Fashion-MNIST 上选得 $T=50$（$\rho_H = 0.021$），在 CIFAR-10 上 $T=5$（$\rho_H = 0.025$）。对高精度模型（MNIST，$\sim97.5\%$），$T=1$ 即可；该选择背后的保序性质在命题~\ref{prop:rank_preserve}（附录~\ref{app:temp}）中证明。对 $T$ 的扫描确认该准则不牺牲信号：在 Fashion-MNIST 与 CIFAR-10 上，softmax RII 在 $T{=}5$ 附近达到峰值，并在 $T{=}10$ 时下降一个数量级以上，随 $T\to\infty$ 趋于零，因此在满足条件数约束的温度中，最小的可接受温度也最可判别。

\subsection{logit 状态与 softmax RII 的桥接}
\label{sec:logit_bridge}

将 logits 量化为 $m$ 个状态（在训练 logits 上做 $k$-means）得到状态分布 $\boldsymbol{\pi}_c\in\Delta^{m-1}$。全文约定 $\chi^2(P\|Q)\coloneqq\sum_i (P_i-Q_i)^2/Q_i$ 表示 $\chi^2$ 散度。状态空间 RII $\rho_{\mathcal{S}}$ 逼近 softmax RII（定理~\ref{thm:state_bridge}，附录~\ref{app:bridge}）；在状态空间中，MHPR 有一个干净的 $\chi^2$ 界：

\begin{theorem}[基于 $\chi^2$ 控制的状态空间 MHPR]
\label{thm:state_mhpr}
设 $\boldsymbol{\pi}_f$ 为遗忘状态分布，$\bar{\boldsymbol{\pi}}_H = \frac{1}{K}\sum_k \boldsymbol{\pi}_{h_k}$。若 $\chi^2(\boldsymbol{\pi}_f \| \bar{\boldsymbol{\pi}}_H) \le \tau^2$ 且 $\|\boldsymbol{\pi}_f\|_2^2 \ge \gamma > 0$，则
\begin{equation}
\rho_{H,\mathcal{S}} \le \frac{\tau^2}{\gamma}.
\label{eq:state_mhpr}
\end{equation}
因此 MHPR 是状态空间输出级不可区分性的一致代理。
\end{theorem}

经验验证（附录~\ref{app:chi2}）确认式~\eqref{eq:state_mhpr} 的界，即对所有梯度上升步均有 $\rho_{H,\mathcal{S}} \le \chi^2/\|\boldsymbol{\pi}_f\|_2^2$。其余状态空间结果（量化误差，引理~\ref{lem:quantization}；谱桥，定理~\ref{thm:state_bridge}；MHPR 偏差率，定理~\ref{thm:softmax_deviation}）推迟到附录~\ref{app:bridge}。

% =============================================================================
% §4 理论结果
% =============================================================================
\section{理论结果}
\label{sec:theory}

\subsection{假设与实践有效性}
\label{sec:assumptions}

表~\ref{tab:assumptions} 列出每个理论结果背后的假设及其在本文实验设定中的状态；随后我们考察其中最具影响的假设。

\begin{table}[t]
\centering
\footnotesize
\caption{每个理论结果背后的假设及其在实验设定中的实践有效性。}\label{tab:assumptions}
\setlength{\tabcolsep}{2pt}
\begin{tabular}{p{2.0cm} p{4.9cm} p{4.8cm}}
\toprule
结果 & 关键假设 & 在我们设定中的有效性 \\
\midrule
定理~\ref{thm:perfect} & $Y$ 是从 softmax 分布 $p_c(x;\theta')$ 抽取的类，而非 argmax 标签 & 按构造成立：所有量都由 softmax 概率向量计算（第~\ref{sec:channel}~节） \\
命题~\ref{prop:confidence} & softmax 向量 sub-Gaussian，方差代理 $\nu^2$；$\sigma_1>0$ & $p_c\in[0,1]$ 有界，故 sub-Gaussian 且 $\nu=1$（下文）；$\sigma_1>0$ 因行是非零概率向量 \\
定理~\ref{thm:finetune} & $L$-Lipschitz 梯度与 softmax；小步条件 $\lambda\|\mathbf{g}\|_2=O(1)$ & 实践中（尤其 LoRA 下，参数空间小）有界；$\lambda{\sim}10^{-5}$ 满足条件 \\
定理~\ref{thm:mi_bound} & $\pi_{\min}=\min_y(\boldsymbol{\mu}_r)_y>0$；等范数（渐近放宽） & 经验通道有 $\pi_{\min}\ge 1/N_r>0$（下文）；等范数渐近放宽（证明见附录~\ref{sec:appendix}） \\
定理~\ref{thm:unified} & 高斯信号模型 $R=\alpha X+\beta Y+N$；小 $\alpha,\eta$ & 理想化；模型下精确，否则一阶成立（注记~\ref{rem:unified_status}） \\
MHPR 命题 & 留出子空间良态（$\sigma_{\min}(\mathbf{H}^*)>0$）；有限样本集中 & 当子空间否则会坍缩时，用 $K\ge2$ 与温度缩放保证（第~\ref{sec:mhpr}~节） \\
\bottomrule
\end{tabular}
\end{table}

\textbf{实践有效性。}(i) \emph{sub-Gaussian softmax。}softmax 向量的每个坐标都落在 $[0,1]$ 内，故有界因而 sub-Gaussian：由 Hoeffding 引理，支撑在 $[0,1]$ 上的变量以 $1/4$ 为方差代理，我们用 $\nu=1$ 作为保守上界。因此命题~\ref{prop:confidence} 以显式常数成立且无未知分布参数。(ii) \emph{$\pi_{\min}>0$。}在每一个实验中每个类在保留集上至少出现一次，故经验 $\pi_{\min}\ge 1/N_r>0$；对类数很多的人群通道，注记~\ref{rem:pi_min} 的截断支撑论证适用。(iii) \emph{$L$-Lipschitz 梯度与 softmax。}softmax 映射对 logit 输入是 $1$-Lipschitz 的，梯度 Lipschitz 常数是损失 Hessian 的谱范数，在本文考虑的（小）参数空间下——尤其 LoRA 适配器只更新少数低秩因子——有界。我们不依赖具体值：定理~\ref{thm:finetune} 的界是小步长时 $\rho=O(\lambda^2)$ 的定性陈述。(iv) \emph{类级遗忘下的交叉协方差。}引理~\ref{lem:cross_cov}（附录~\ref{app:finetune}）依赖不同类别的 logit 梯度在良训练分类器下近似正交，这在分类器自信（$p_f(c_f)\approx1$）时成立；附录给出推导。

\subsection{伪代码}
\label{sec:pseudocode}

为可复现性，我们以伪代码给出 RII、MHPR 与推荐审计协议的计算（伪代码与英文版一致，保持代码为英文）；实现见 \url{https://github.com/EurusDemerzel/rii_unlearning}。

{\small\begin{verbatim}
Algorithm 1: Compute RII rho
  Input:  model theta', forget set D_f, retain set D_r
  1. mu_f <- mean_x in D_f Softmax(f_{theta'}(x))
  2. mu_r <- mean_x in D_r Softmax(f_{theta'}(x))
  3. M <- [mu_f ; mu_r]            # 2 x C channel matrix
  4. sigma_1, sigma_2 <- singular values of M
  5. rho <- 1 - sigma_1^2/(sigma_1^2 + sigma_2^2)
  Output: rho

Algorithm 2: Compute MHPR rho_H
  Input:  model theta', forget set D_f, held-out means H = [mu_h1 ... mu_hK]
  1. mu_f <- mean_x in D_f Softmax(f_{theta'}(x))
  2. P <- H^+ H                     # projection onto held-out span
  3. rho_H <- ||mu_f - P mu_f||^2 / ||mu_f||^2
  Output: rho_H

Algorithm 3: Recommended audit protocol
  1. Sample N_min examples from D_f, D_r (and each held-out class for MHPR)
  2. Compute rho (or rho_H); apply temperature scaling if accuracy <~ 90%
  3. Report [rho - eps, rho + eps] as a high-probability interval (1 - delta)
  4. Report MIA AUC alongside RII for a two-axis assessment
\end{verbatim}}

下面的结果关注\emph{通道}——类别输出 $Y$ 的条件发射分布——而非 softmax 向量的逐样本分布；这一区分的范围在注记~\ref{rem:channel} 中讨论。

\begin{theorem}[输出分布不可逆性]
\label{thm:perfect}
设 $X\in\{0,1\}$ 为遗忘集成员指示变量（$X{=}1$ 表示遗忘、$X{=}0$ 表示保留，下文简写为 $X{=}f$ 与 $X{=}r$），$Y\in\{1,\dots,C\}$ 为模型在 $\theta'$ 下的类别输出：对输入 $x$，通道以概率 $p_c(x;\theta')$ 发射 $Y{=}c$，从而由全概率公式 $P(Y{=}c\mid X{=}f)=(\boldsymbol{\mu}_f)_c$ 且 $P(Y{=}c\mid X{=}r)=(\boldsymbol{\mu}_r)_c$（注记~\ref{rem:channel}）。则
\begin{equation}
\rho(\theta') = 0 \;\Longleftrightarrow\; \boldsymbol{\mu}_f = \boldsymbol{\mu}_r \;\Longleftrightarrow\; I(X; Y \mid \theta') = 0.
\end{equation}
因此 $\rho=0$ 是\emph{通道级}（均值输出）分布不可逆性的充要条件：通道的类别发射独立于遗忘集成员。逐样本输出分布不由 $\rho$ 证明（注记~\ref{rem:channel}）。
\end{theorem}

\begin{proof}
$\rho=0 \Rightarrow \sigma_2=0 \Rightarrow \mathbf{M}$ 秩为 1。两行都是和为 1 的概率向量，因此 $\boldsymbol{\mu}_f = \lambda\boldsymbol{\mu}_r$ 且 $\lambda>0$；和为一迫使 $\lambda=1$，故 $\boldsymbol{\mu}_f=\boldsymbol{\mu}_r$。由全概率公式，$(\boldsymbol{\mu}_f)_c = \frac{1}{N_f}\sum_{x\in D_f}p_c(x;\theta') = \frac{1}{N_f}\sum_{x\in D_f}P(Y{=}c\mid x;\theta') = P(Y{=}c\mid X{=}f)$，同理 $(\boldsymbol{\mu}_r)_c = P(Y{=}c\mid X{=}r)$；因此 $\boldsymbol{\mu}_f=\boldsymbol{\mu}_r$ 恰为 $P(Y\mid X{=}f)=P(Y\mid X{=}r)$，即 $Y\perp X\mid\theta'$ 且 $I(X;Y\mid\theta')=0$。反向立即成立：$I(X;Y\mid\theta')=0$ 迫使 $P(Y\mid X{=}f)=P(Y\mid X{=}r)$，即 $\boldsymbol{\mu}_f=\boldsymbol{\mu}_r$，故 $\sigma_2=0$ 且 $\rho=0$。
\end{proof}

\begin{remark}[该准则证明什么——以及不证明什么]
\label{rem:channel}
一个自然的顾虑是 $\boldsymbol{\mu}_f=\boldsymbol{\mu}_r$ 只是一阶矩条件：逐样本输出\emph{均值}相等并不意味着 $\mathbf{p}(x)$ 在 $D_f$ 与 $D_r$ 上的\emph{分布}相等。此顾虑针对的是\emph{逐样本}不可区分性，而非\emph{通道}证书。因为 $Y$ 是取值 $C$ 的类别变量，其条件分布 $P(Y\mid X{=}f)$ \emph{就是} $C$ 维向量 $\boldsymbol{\mu}_f$（由全概率公式），所以均值相等恰等于通道条件发射分布相等：无需指数族或充分统计量假设。若 $Y$ 是连续 softmax 向量本身，一阶矩保留警告才会生效；那是逐样本层面，$\rho$ 不证明它，而由 PSSD 得分（定义~\ref{def:isd}，附录~\ref{app:pssd}）处理：两个集合可以共享相同通道，同时个别样本保留异常的 logit 状态，因此即使 $\rho=0$，PSSD 也可能非零。
\end{remark}

\begin{proposition}[有限样本置信]
\label{prop:confidence}
设 $\rho^*$ 为真实通道矩阵 $\mathbf{M}^*$ 的总体 RII，$\hat{\rho}$ 为由 $N_f, N_r$ 个样本得到的估计。对方差代理为 $\nu^2$ 的次高斯 softmax 向量，
\begin{equation}
\mathbb{P}\!\left(|\hat{\rho} - \rho^*| \le \frac{4\nu}{\sigma_1}\sqrt{\frac{C\log(4C/\delta)}{N_{\min}}}\right) \ge 1-\delta,
\end{equation}
其中 $N_{\min}=\min(N_f,N_r)$。要以置信度 $1-\delta$ 实现误差 $\le\varepsilon$，需从每个集合采样 $N_{\min} \ge 16C\nu^2\log(4C/\delta)/(\sigma_1^2\varepsilon^2)$。
\end{proposition}

\begin{theorem}[微调泄漏界]
\label{thm:finetune}
对单步梯度上升 $\theta' = \theta_0 + \lambda\mathbf{g}$（$\mathbf{g}=\nabla_\theta\mathcal{L}(\theta_0, D_f)$），其中梯度 $L$-Lipschitz 且 softmax 映射 $L$-Lipschitz，且步长 $\lambda$ 满足小步条件 $\lambda\|\mathbf{g}\|_2=O(1)$，
\begin{equation}
\rho(\theta') \le \frac{\lambda^2 L^2}{2\sigma_1^2}\,\operatorname{Tr}\bigl(\boldsymbol{\Sigma}_f\bigr) + O(\lambda^3),
\label{eq:finetune_bound}
\end{equation}
其中 $\boldsymbol{\Sigma}_f \coloneqq \operatorname{Cov}_{x\sim D_f}(\nabla_\theta f_{\theta_0}(x))$ 是遗忘集上 logit 映射的梯度协方差，$\sigma_1$ 是 $\theta_0$ 处通道的首个奇异值。不假设对 $\mathbf{g}$ 做归一化；其尺度通过小步条件与 $\operatorname{Tr}(\boldsymbol{\Sigma}_f)$（它下界平均平方梯度范数：$\operatorname{Tr}(\boldsymbol{\Sigma}_f)=\mathbb{E}\|\nabla_\theta f-\mathbb{E}\nabla_\theta f\|_2^2\le\mathbb{E}\|\nabla_\theta f\|_2^2$）进入。因此小步长时 $\rho = O(\lambda^2)$。
\end{theorem}

\begin{remark}[定理~\ref{thm:finetune} 的适用范围]
该界是局部的：它在轨迹停留在 $\theta_0$ 的 $L$-Lipschitz 邻域内时成立，这由小步条件保证。等价地，可将上升方向归一化并把 $1/\|\mathbf{g}\|_2$ 吸收进 $\lambda$ 的定义；除有效步长的重新缩放外，陈述不变。在实际使用的学习率（$\lambda{\sim}10^{-5}$，$\|\mathbf{g}\|_2$ 量级为 1）下，该条件被舒适地满足。
\end{remark}

\paragraph{类级遗忘：受控误差分析。}
对类级遗忘（$D_f$ = 类 $c_f$ 的全部样本），$D_f$ 与 $D_r$ 上的梯度相关。交叉熵损失有特殊结构，使我们能控制这种相关。其 logit 梯度为 $\nabla_z\ell(z,y)=\mathbf{p}(z)-\mathbf{e}_y$，其中 $\mathbf{e}_y$ 为 one-hot 向量。对来自不同类 $c_f\neq c_r$ 的样本：
\begin{equation}
\langle\nabla_z\ell_f,\nabla_z\ell_r\rangle = \langle\mathbf{p}_f,\mathbf{p}_r\rangle - p_f(c_r) - p_r(c_f).
\end{equation}
在 $\theta_0$（训练良好的分类器）处，$p_f(c_f){\approx}1$ 且 $p_r(c_r){\approx}1$，故此内积 $\approx0$。经过 $\lambda$ 梯度步后，logits 位移 $O(\lambda)$，且 softmax $1$-Lipschitz 连续性给出 $\|\nabla_z\ell' - \nabla_z\ell\|_2\le\|\Delta z\|_2 = O(\lambda)$。因此 $|\langle\nabla_z\ell'_f,\nabla_z\ell'_r\rangle| = O(\lambda)$，这意味着参数空间交叉协方差 $\boldsymbol{\Sigma}_{fr} \coloneqq \operatorname{Cov}(\nabla_\theta\ell_f,\nabla_\theta\ell_r)$ 满足 $\|\boldsymbol{\Sigma}_{fr}\|_F = O(\lambda)$（引理~\ref{lem:cross_cov}，附录~\ref{app:finetune}）。

因此，类级界保持式~\eqref{eq:finetune_bound} 的 $O(\lambda^2)$ 阶：其 $O(\lambda^3)$ 项吸收交叉协方差贡献（附录~\ref{app:finetune}），且由于 $\|\boldsymbol{\Sigma}_{fr}\|_F = O(\lambda)$ 只增加高阶修正，在实际学习率（$\lambda{\sim}10^{-5}$）下可忽略，故该阶得以保持。

\begin{theorem}[谱 $\chi^2$--MI 界]
\label{thm:mi_bound}
对 $\theta'$ 下的二元成员 $X$ 与预测类 $Y$，当通道两行等范数 $\|\boldsymbol{\mu}_f\|_2=\|\boldsymbol{\mu}_r\|_2$（该条件在渐近意义下放宽，见下文）时，
\begin{equation}
I(X; Y \mid \theta') \le \frac{\sigma_1^2}{\pi_{\min}}\cdot\frac{\rho}{1-\rho},
\end{equation}
其中 $\pi_{\min} \coloneqq \min_y (\boldsymbol{\mu}_r)_y > 0$ 且 $(\boldsymbol{\mu}_r)_y=P(Y{=}y\mid X{=}r)$。当 $\rho\ll 1$ 时，$I \le O(\rho)$；该界随 $\rho\to0$ 消失，并在无等范数条件时渐近成立（证明见附录~\ref{sec:appendix}）。
\end{theorem}

\begin{remark}[关于 $\pi_{\min}>0$]
\label{rem:pi_min}
对经验通道，$(\boldsymbol{\mu}_r)_y$ 是类 $y$ 在保留集上的相对频率，因此只要每个类至少出现一次（本文所有实验的情形）就有 $\pi_{\min}>0$。对类数很多、某些类概率任意小的人群通道，可将该界应用于截断支撑 $\mathcal{Y}_\epsilon=\{y:(\boldsymbol{\mu}_r)_y>\epsilon\}$，代价是 $O(1/\epsilon)$ 的乘法因子；同样的截断出现在置信界中（命题~\ref{prop:confidence}）。
\end{remark}

\paragraph{信号模型。}
我们将观测到的移除后输出建模为 $R = \alpha X + \beta Y + N$，$N \sim \mathcal{N}(0,\sigma^2)$，其中 $X$ 为遗忘/保留指示，$Y$ 为通道输出，$\alpha,\beta$ 为耦合系数。这刻画两条正交的泄漏通路：直接信号泄漏（$\propto \alpha^2$，由 MIA 检测）与覆盖泄漏（$\propto \eta^2$，即对秩一的谱偏离，由 RII 检测）。

\begin{lemma}[小 $\alpha$ Fisher 信息展开]
\label{lem:fisher}
对 $R = \alpha X + Z$（$X \perp Z$），当 $\alpha \to 0$ 时：
$I(X;R) = \frac{\alpha^2}{2}\mathrm{Var}(X)J(Z) + o(\alpha^2)$，其中 $J(Z)$ 为 $Z$ 的 Fisher 信息。对高斯噪声，$J(N) = 1/\sigma^2$。
\end{lemma}

\begin{theorem}[统一泄漏分解]
\label{thm:unified}
在信号模型 $R = \alpha X + \beta Y + N$ 与通道奇异值 $\{\sigma_k\}$ 下，
\begin{equation}
I(X; R) \le \frac{\alpha^2}{2\sigma^2}\operatorname{Var}(X) \;+\; \eta^2 \;+\; o(\alpha^2 + \eta^2),
\end{equation}
其中 $I_{\mathrm{dir}} \propto \alpha^2$（直接泄漏，由 MIA 测量）且 $I_{\mathrm{cov}} = \eta^2 = \sum_{k\ge 2}\sigma_k^2$（覆盖泄漏，由 RII 测量）。
\end{theorem}

\begin{remark}[定理~\ref{thm:unified} 的地位]
\label{rem:unified_status}
该分解在高斯信号模型下精确成立，并在引理~\ref{lem:fisher} 的正则性下对 $\alpha^2+\eta^2$ 一阶成立（小 $\alpha$、小谱偏离）；为可读性我们以定理陈述之，证明草图见附录~\ref{app:unified}，完整的 $\chi^2$ 控制在附录~\ref{app:bridge_proofs}。
\end{remark}

\paragraph{正交泄漏轴。}
MIA 准确率跟踪直接轴 $I_{\mathrm{dir}}$（分布位移）；RII 跟踪覆盖轴 $I_{\mathrm{cov}}$（输出级不可区分性）。两者可以同时很大：它们是泄漏的独立维度（见第~\ref{sec:mia}~节）。

\paragraph{连续类比。}
离散 $2\times C$ 分析通过 Hilbert--Schmidt 算子理论提升到连续空间，其中偏离参数 $\eta^2=\sum_{k\ge2}\sigma_k^2$ 是定理~\ref{thm:unified} 的基础（附录~\ref{app:bridge}）。

% =============================================================================
% zh_part2.tex — §5 实验评估（逐句对照中文版）
% =============================================================================

% =============================================================================
\section{实验评估}
\label{sec:experiments}

\subsection{设置}

\textbf{数据集与模型：}MNIST（60K，10 类）、Fashion-MNIST（60K，10）、CIFAR-10（50K，10）、CIFAR-100（50K，100；类级协议使用 20 类训练子集）。架构：SimpleMLP（784$\to$128$\to$10）、SmallCNN（3 卷积 + 2 全连接）、ResNet-18、DeepMLP（用于过拟合）。训练：Adam，lr $10^{-3}$，batch 64，10 epoch。语言模型研究（第~\ref{sec:llm}~节）使用 LLaMA-2-7B 上的 TOFU~\cite{maini2024tofu} 作者级基准，采用基于 LoRA 的 PEFT。

\textbf{遗忘方法：}NoUnlearn（不移除）、Retrain（黄金标准）、SISA~\cite{bourtoule2021machine}（$S{=}5$ 分片）、NegGrad（仅在 $D_f$ 上梯度上升）、FineTune（在 $D_f$ 上梯度上升后接在 $D_r$ 上下降）、KED（通过 KL 到均匀的知识侵蚀，作用于 $D_f$）、以及 BadTeacher（在 $D_f$ 上进行错误标签训练）。

\textbf{基准协议。}在第~\ref{sec:benchmark}~节中，我们在 CIFAR-10 上采用标准类级协议：模型在类 $\{0,\dots,6\}$ 上训练，类~3（cat）为遗忘集，类 $\{7,8,9\}$ 留出作为未见 MHPR 参照（$K{=}3$）。Fashion-MNIST 使用相同协议：七个训练类 $\{0,\dots,6\}$，遗忘类~3（Dress），留出类 $\{7,8,9\}$（Sneaker、Bag、Ankle boot）。两种情形下的留出类都是训练切分后的剩余类——按可用性而非语义相似性选择；该选择的影响在第~\ref{sec:benchmark}~节（注意事项~(2)）讨论。为确定性，模型在无数据增强下训练，且每个遗忘方法都用\emph{相同}的指标组评估。

在全文范围内，``MIA''随协议而异，指代\emph{三种不同}的度量；为避免歧义，每个表格都标明其所指：随机遗忘基线中的训练/测试成员推断准确率（表~\ref{tab:main}）、类级基准中的遗忘/保留损失阈值 AUC（表~\ref{tab:benchmark}）、以及 TOFU 上的遗忘集平均下一词 NLL（表~\ref{tab:llm}）。在对比两条泄漏轴时（第~\ref{sec:mia}~节），``MIA''指训练/测试成员准确率。

跨下面的实验，标准训练产生近秩一输出通道（$\rho<5\times10^{-3}$），使 RII 成为稳定、无需调参的量：随机样本级遗忘平凡地不可区分，而类级遗忘不是（第~\ref{sec:benchmark}~节）。在类级遗忘下，RII 与 MHPR 对遗忘方法的排序与遗忘强度一致，并揭示出留下（甚至增强）输出签名而准确率与 MIA 都错过的那些方法；在 TOFU 基准上的 70 亿参数语言模型上（第~\ref{sec:llm}~节），即使所有已遗忘检查点都将生成准确率饱和在 100\%，同样排序仍然成立。最后，RII 与 MIA 度量重叠但不同的泄漏轴（$r{=}0.57$），因此完整审计应同时报告两者。

\subsection{基线：标准训练下的近秩一输出}

虽然本工作聚焦类级遗忘，我们首先建立一个简短的随机（样本级）基线：它分离出标准训练的近秩一结构，并表明 RII 并不退化，因为随机子集平凡地不可区分，而内聚的语义类（第~\ref{sec:mhpr}~节与第~\ref{sec:benchmark}~节）则不然。这些结果仅作为控制基线。

\begin{table}[t]
\centering
\small
\caption{秩一基线：标准训练下的 RII（5\% 随机遗忘；所有行均为 NoUnlearn）。MIA 是训练/测试成员推断准确率（第~\ref{sec:mia}~节），而非遗忘/保留损失阈值 AUC。}\label{tab:main}
\setlength{\tabcolsep}{2pt}
\begin{tabular}{l c c c c}
\toprule
数据集 & 方法 & Acc.(\%) & $\rho$ & MIA(\%) \\
\midrule
MNIST & NoUnlearn & 97.45 & $9.762\times10^{-4}$ & 90.4 \\
\midrule
CIFAR-10 & NoUnlearn & 76.1 & $1.20\times10^{-3}$ & 90.5 \\
CIFAR-100 (SmallCNN) & NoUnlearn & 20.8 & $7.10\times10^{-4}$ & 90.5 \\
Fashion-MNIST & NoUnlearn & 88.7 & $2.83\times10^{-4}$ & 90.4 \\
\bottomrule
\end{tabular}
\end{table}

表~\ref{tab:main} 揭示四点关键发现。第一，$\rho<5\times10^{-3}$ 一致成立：标准神经网络训练产生近秩一输出通道。第二，$\rho$ 与准确率无关：CIFAR-100 在 20\% 准确率下给出的 $\rho$ \emph{低于} MNIST 在 97.5\% 下的值，证实秩一性质是 softmax 输出单纯形固有的。第三，不同遗忘方法给出几乎相同的 $\rho$（我们在 MNIST 上验证了 Retrain、SISA 与 FineTune），这是预期的：对随机子集，模型的泛化使输出即使不做显式遗忘也不可区分。第四，在 10 个随机种子下，$\rho$ 的标准差 $<3.5\times10^{-4}$，支持可靠的单次审计。因此在随机遗忘下 RII 近乎平凡（所有方法不可区分），而类级遗忘是 RII 与 MHPR 变得可判别的区间（第~\ref{sec:mhpr}~节与第~\ref{sec:benchmark}~节）。

\subsection{70 亿参数语言模型上的 RII（TOFU）}
\label{sec:llm}

为检验该证书能否经受尺度与模态的迁移，我们接下来转向当代指令微调语言模型。我们在 \emph{LLaMA-2-7B}（官方 TOFU 微调检查点）上运行标准 TOFU~\cite{maini2024tofu} 作者级协议，遗忘集 $=$ 一位作者（40 个 QA 对），保留划分为 90 位作者，采用实践者实际部署的 PEFT 方案：LoRA（$r{=}8$，scale~20，作用于 32 层中的 8 层）。方法：NoUnlearn（检查点原样）、FineTune（在保留答案上 60 步 LoRA，AdamW $10^{-5}$）、NegGrad（在遗忘答案上 10 步 LoRA 梯度上升，SGD $2{\times}10^{-5}$），以及 Retrain oracle（基座 LLaMA-2-7B 加仅在保留集上的 LoRA；基座模型从未见过遗忘作者）。RII 在遗忘与保留 QA 对的平均下一 token softmax 分布构成的 $2\times V$（$V{=}32{,}000$）通道矩阵上计算；MIA 为遗忘集的平均下一 token NLL。

\begin{table}[t]
\centering
\small
\caption{LLaMA-2-7B 上 TOFU 基准的作者级遗忘（基于 LoRA；Retrain 为轻量 60 步 oracle）。}\label{tab:llm}
\setlength{\tabcolsep}{2pt}
\begin{tabular}{l c c c c}
\toprule
方法 & Forget(\%) & Retain(\%) & RII & MIA \\
\midrule
NoUnlearn & 100.0 & 100.0 & $1.16{\times}10^{-1}$ & 1.348 \\
FineTune & 100.0 & 100.0 & $9.73{\times}10^{-2}$ & 0.852 \\
NegGrad & 100.0 & 100.0 & $\mathbf{1.17{\times}10^{-1}}$ & 1.352 \\
Retrain (oracle) & 80.0 & 65.0 & $\mathbf{8.01{\times}10^{-2}}$ & 2.163 \\
\bottomrule
\end{tabular}
\par\vspace{2pt}\footnotesize MIA 是遗忘集的平均下一 token NLL，\emph{不是} AUC；值越大表示遗忘越强。
\end{table}

三个遗忘后的检查点（以及未触动模型）都能完美回答遗忘问题（Forget$=$100\%；许多 TOFU 答案可以直接从问题本身恢复），因此生成准确率无法区分它们。相比之下，RII 将 FineTune 排在 NoUnlearn 之下（$9.73{\times}10^{-2}$ 对比 $1.16{\times}10^{-1}$，$-16\%$），与视觉基准的排序相同，并将（轻量、60 步的）重训练 oracle（$8.01{\times}10^{-2}$）与未触动模型区分开 $1.45\times$。oracle 另外还可由生成准确率（80\% 对比 100\%）与其大得多的 MIA（2.163 对比 1.348）区分；此处的声明更窄：在生成准确率饱和于 100\% 的检查点中，RII 提供了准确率无法给出的排序。我们不声称仅凭 RII 就能区分 NegGrad 与 NoUnlearn：它们的全集差距（$1.7{\times}10^{-3}$）落在评估噪声内，且两者子采样均值一致（均为 $0.127\pm0.002$，见下文）。

NegGrad 的 RII（$1.17{\times}10^{-1}$）在我们的评估协议下与 NoUnlearn（$1.16{\times}10^{-1}$）统计上不可区分，不同于 FineTune 的 16\% 下降；其 MIA（1.352）同样匹配 NoUnlearn（1.348）。因此上升在此不产生可测量的输出空间擦除，而下面的强度扫描表明，随着上升加强，遗忘输出签名\emph{增强}（在 $3\times$ 强度时升至 $1.20{\times}10^{-1}$，随后模型崩溃），这是第~\ref{sec:benchmark}~节 CIFAR-10/100 上所述隐藏失败的 7B 类比。

在保留答案上继续 LoRA 训练使遗忘与保留输出分布靠拢：RII 降到 $9.73{\times}10^{-2}$，下降 $16\%$。它还取得最低 MIA（0.852），意味着模型对遗忘集仍最流畅。这与视觉结果一致。然而，如第~\ref{sec:rebound}~节所述，RII 降低并非真擦除的充分条件，因为遗忘准确率仍为 100\%。

为衡量评估稳定性，我们对评估集做三次随机子采样重复评估；每种方法的 RII 标准差低于 $3.5{\times}10^{-3}$（相对 $<$3\%），且排序 Retrain $<$ \allowbreak FineTune $<$ \allowbreak NoUnlearn ${\approx}$ \allowbreak NegGrad 在所有种子下不变，NegGrad 与 NoUnlearn 之间的差距（$1.7{\times}10^{-3}$）在噪声内。子采样均值（基于子采样集计算，因此不与表~\ref{tab:llm} 的完整集数值直接可比）分别为 $0.088\pm0.003$（oracle）、$0.107\pm0.002$（FineTune），NoUnlearn 与 NegGrad 均为 $0.127\pm0.002$。

在小模型 Qwen2.5-0.5B pilot（第~\ref{sec:benchmark}~节）上，词表级通道无法将 oracle 与未触动模型区分开，因为在 0.5B 规模下两个答案集诱导出统计上相似的文本分布。在 7B 规模下，区分显著（$1.45\times$，三次子采样区间无重叠）：更大模型在下一 token 分布中更强地编码作者身份，因此即使原始词表通道也携带类级信号。因此通道\emph{语义}决定判别力（类级头仍是推荐的审计对象），而模型容量决定原始通道是否足够（参见注记~\ref{rem:channel}）。

我们还对梯度上升强度（步数 $\times$ 学习率）做扫描，从基线 $10\times 2{\times}10^{-5}$ 到 $30\times 10^{-4}$（$3\times$ 步数，$5\times$ lr）与 $60\times 2{\times}10^{-4}$（$6\times$ 步数，$10\times$ lr），每次都从全新的零 LoRA 出发。在中等强度下，RII 进一步\emph{上升}到 $1.20{\times}10^{-1}$（高于 NoUnlearn 的 $1.16{\times}10^{-1}$），而遗忘/保留准确率保持在 100\%：隐藏失败并不饱和——更强的上升锐化遗忘输出签名。在最强设置下模型开始崩溃：MIA 跳到 5.25（从 1.35），上升损失发散到 6.5，而 RII 回落到 $1.01{\times}10^{-1}$，但生成准确率仍是 100\%，因此仅凭准确率仍无法检测崩溃。此处 RII 下降\emph{是因为输出分布退化为近随机噪声}，这巧合地使两个集合看起来又统计相似；崩溃反而由 MIA 标记。因此 RII 是双向的：它随隐藏失败签名增强而上升，随输出分布本身退化而下降，所以低 RII 必须与效用指标（保留准确率、MIA）一起解读，以区分真擦除与模型崩溃，正如第~\ref{sec:mia}~节的两轴审计所建议。

\subsection{统一协议下的基准比较}
\label{sec:benchmark}

表~\ref{tab:benchmark} 在上述协议下用单一指标组评估七种遗忘方法：保留/遗忘准确率、RII $\rho$、MHPR（$K{=}3$）、MIA（损失阈值 AUC）、倒数第二层特征上的表示级 MMD，以及试图分离遗忘类特征与留出未见类特征的残差线性探针。

\begin{table}[t]
\centering
\small
\caption{CIFAR-10 上的类级遗忘基准（遗忘 cat，$K{=}3$ 留出）；三个种子的均值 $\pm$ 标准差。}\label{tab:benchmark}
\setlength{\tabcolsep}{3pt}
\begin{tabular}{l c c c c c}
\toprule
方法 & Ret.(\%) & For.(\%) & $\rho$ & MHPR & MIA-loss \\
\midrule
NoUnlearn & $92.6{\pm}0.5$ & $85.2{\pm}10.7$ & $0.190{\pm}0.036$ & $0.802{\pm}0.023$ & $0.306{\pm}0.077$ \\
Retrain (oracle) & $94.5{\pm}0.9$ & 0.0 & $0.113{\pm}0.005$ & $0.468{\pm}0.115$ & 0.000 \\
NegGrad & $91.9{\pm}1.5$ & $63.8{\pm}10.0$ & $\mathbf{0.249{\pm}0.002}$ & $0.839{\pm}0.046$ & $0.168{\pm}0.040$ \\
FineTune & $98.0{\pm}0.3$ & $28.5{\pm}3.9$ & $0.212{\pm}0.008$ & $0.668{\pm}0.053$ & $0.029{\pm}0.004$ \\
KED & $97.3{\pm}0.3$ & $1.2{\pm}0.4$ & $0.134{\pm}0.009$ & $0.459{\pm}0.035$ & $0.003{\pm}0.001$ \\
BadTeacher & $97.7{\pm}0.3$ & $13.2{\pm}5.9$ & $0.161{\pm}0.015$ & $0.367{\pm}0.058$ & $0.011{\pm}0.006$ \\
SISA & $70.4{\pm}0.1$ & 0.0 & $0.073{\pm}0.004$ & $0.077{\pm}0.007$ & $0.002{\pm}0.000$ \\
\bottomrule
\end{tabular}
\par\vspace{2pt}\footnotesize 三个种子的均值 $\pm$ 标准差。MIA-loss 是损失阈值 AUC；Retrain 的 $0.000$ 是预期的（oracle 从未在遗忘类上训练，产生完美反向排序；见第~\ref{sec:benchmark}~节）。表示 MMD 与残差探针 AUC 在正文中报告（第~\ref{sec:benchmark}~节）。
\end{table}

即使是重训练 oracle（数学上完美的遗忘）也有非零 RII，$\rho{=}0.113{\pm}0.005$：因为 cat 现在是\emph{未见}类，其输出分布不同于保留类，所以标准 RII 无法在类级遗忘下证明 oracle 完美。这种已知/未知不对称正是 MHPR 所处理的，它将与 NoUnlearn 的差距缩小到（$0.468$ 对比 $0.802$）。

梯度上升（NegGrad）是\emph{唯一} RII 随遗忘\emph{上升}的方法（$0.249{\pm}0.002$ 对比 NoUnlearn 的 $0.190{\pm}0.036$，表~\ref{tab:benchmark}），尽管遗忘准确率与 MIA-loss 都改善了。上升将遗忘类推向自信、系统性错误的预测，集中输出签名并使通道\emph{更加}可区分；在我们报告的指标中，RII 是唯一暴露此失败的，因为准确率与 MIA-loss 都改善。结合微调反弹（第~\ref{sec:rebound}~节），输出空间擦除在基于梯度的遗忘强度上不是单调的。

\begin{table}[t]
\centering
\small
\caption{七种遗忘方法间评估指标的 Pearson 相关（CIFAR-10；每个方法是三个种子的均值，七个方法级数据点仅供参考）。}
\label{tab:corr}
\setlength{\tabcolsep}{4pt}
\begin{tabular}{l c c c c c c}
\toprule
 & RII & MHPR & For. & MIA & MMD & probe \\
\midrule
RII & 1.00 & & & & & \\
MHPR & 0.90 & 1.00 & & & & \\
For. & 0.75 & 0.83 & 1.00 & & & \\
MIA & 0.57 & 0.73 & 0.96 & 1.00 & & \\
MMD & 0.77 & 0.69 & 0.58 & 0.47 & 1.00 & \\
probe & 0.80 & 0.75 & 0.48 & 0.35 & 0.95 & 1.00 \\
\bottomrule
\end{tabular}
\end{table}

RII 与 MHPR 强一致（$r{=}0.90$）并跟踪遗忘准确率（$r{=}0.75$ 与 $0.83$）；RII--遗忘准确率相关由 Retrain 端点（遗忘准确率 $0\%$）驱动，且在非 oracle 方法中这两个指标系统性地分歧：NegGrad 的遗忘准确率\emph{下降}（从 $85.2\%$ 到 $63.8\%$）而 RII\emph{上升}（从 $0.190$ 到 $0.249$），这是准确率代理会错过的失败。RII 与 MIA-loss 呈中等相关（$r{=}0.57$），并在顶端分歧（NegGrad 在 RII 下最差、在 MIA-loss 下第二差），支持正交轴观点（第~\ref{sec:mia}~节）。表示 MMD 与 RII 相关 $r{=}0.77$；完整矩阵见表~\ref{tab:corr}。

分离遗忘类与留出特征的残差线性探针在\emph{所有}方法上都成功（$\mathrm{AUC}\approx0.94$--$0.98$），包括 oracle（$0.956$）：输出级擦除与表示级签名共存，与~\cite{cosma2026ruler,yong2026erased} 一致，并促使两层审计（第~\ref{sec:mia}~节）。

解读这些结果需要若干注意事项。(1) SISA 的低保留准确率（$70.4\%$）源于轻量训练预算（$5$ 分片 $\times$ $2$ 切片），因此其近零谱得分是效用伪影，我们将其排除出排序结论：退化的模型可能把 RII 推向零（如第~\ref{sec:llm}~节与第~\ref{sec:benchmark}~节中的崩溃），所以低 RII 应与保留准确率一起解读。(2) CIFAR-10 上 MHPR 绝对值高于 MNIST，因为 $\{7,8,9\}$ 与 cat 语义上不近；当留出参照有代表性时 MHPR 最具判别力。(3) MIA-loss AUC 低于 $0.5$，因为 cat 比保留类更难（遗忘/保留准确率 $85.2\%$ 对比 $92.6\%$），因此损失阈值沿``错误''方向分离它们；相对排序是有意义的。Retrain oracle 的极端值 $0.000$ 是预期的而非异常：因为 oracle 从未在 cat 上训练，其 cat 损失超过每个保留样本，所以基于损失的 MIA 将每个遗忘样本判给非成员类，即完美反向排序（AUC${\to}0$），这本身就是成功擦除的分布级证据。(4) 以下类级基准基于三个种子，以均值 $\pm$ 标准差报告；它们支持指标的稳定性与定性排序，但在 $n{=}3$ 下并不单独确立相近对（例如 RII 下 NegGrad 与 NoUnlearn）的统计显著性。RII 的种子稳定性还由随机基线（10 种子，SD$<3.5{\times}10^{-4}$）与 7B 研究（3 次子采样，相对 $<$3\%）交叉佐证。

表~\ref{tab:cross} 用 MLP 在 Fashion-MNIST 上重复该协议。RII 仍可判别：NegGrad 失败持续（$\rho{=}0.169{\pm}0.004$ 对比 $0.166{\pm}0.004$），oracle 保持非零 RII（$0.089{\pm}0.012$），KED 取得最低 RII（$0.063{\pm}0.006$）。与早前分析的单次基线不同，MHPR 在全部三个种子上不加温度缩放即可判别：oracle（$0.429{\pm}0.137$）排在 NoUnlearn（$0.749{\pm}0.244$）\emph{之下}，而在单次较早运行中观察到的低精度子空间坍缩（第~\ref{sec:mhpr}~节）是种子依赖的而非系统性的；温度缩放仍是一般情形下低精度模型的推荐保障。表示探针再次在 $\mathrm{AUC}{=}1$ 附近饱和。

\begin{table}[t]
\centering
\small
\caption{相同协议下的跨数据集基准：CIFAR-10（SmallCNN）与 Fashion-MNIST（MLP）上的 RII/MHPR/MIA；三个种子的均值 $\pm$ 标准差。}\label{tab:cross}
\setlength{\tabcolsep}{2pt}
\footnotesize
\begin{tabular}{l c c c | c c c}
\toprule
 & \multicolumn{3}{c|}{CIFAR-10} & \multicolumn{3}{c}{Fashion-MNIST} \\
方法 & $\rho$ & MHPR & MIA & $\rho$ & MHPR & MIA \\
\midrule
NoUnlearn & $0.190{\pm}0.036$ & $0.802{\pm}0.023$ & $0.306{\pm}0.077$ & $0.166{\pm}0.004$ & $0.749{\pm}0.244$ & $0.493{\pm}0.046$ \\
Retrain (oracle) & $0.113{\pm}0.005$ & $0.468{\pm}0.115$ & 0.000 & $0.089{\pm}0.012$ & $0.429{\pm}0.137$ & 0.000 \\
NegGrad & $\mathbf{0.249{\pm}0.002}$ & $0.839{\pm}0.046$ & $0.168{\pm}0.040$ & $0.169{\pm}0.004$ & $0.787{\pm}0.197$ & $0.474{\pm}0.049$ \\
FineTune & $0.212{\pm}0.008$ & $0.668{\pm}0.053$ & $0.029{\pm}0.004$ & $0.245{\pm}0.028$ & $0.697{\pm}0.106$ & $0.115{\pm}0.029$ \\
KED & $0.134{\pm}0.009$ & $0.459{\pm}0.035$ & $0.003{\pm}0.001$ & $\mathbf{0.063{\pm}0.006}$ & $0.213{\pm}0.062$ & $0.008{\pm}0.002$ \\
BadTeacher & $0.161{\pm}0.015$ & $0.367{\pm}0.058$ & $0.011{\pm}0.006$ & $0.129{\pm}0.010$ & $0.169{\pm}0.096$ & $0.008{\pm}0.003$ \\
SISA & $0.073{\pm}0.004$ & $0.077{\pm}0.007$ & $0.002{\pm}0.000$ & $0.096{\pm}0.008$ & $0.592{\pm}0.064$ & $0.003{\pm}0.001$ \\
\bottomrule
\end{tabular}
\end{table}

表~\ref{tab:cifar100} 用 20 个训练类重复该协议（遗忘类~3，留出 $\{20,21,22\}$）。NegGrad 隐藏失败在更大规模下\emph{更显著}：其 RII（$0.189{\pm}0.016$）甚至超过 NoUnlearn（$0.127{\pm}0.026$，$1.49\times$ 差距，对比 CIFAR-10 的 $1.31\times$），而遗忘准确率与 MIA-loss 再次改善：被梯度上升增强的输出签名在更大的输出空间中更尖锐。已知/未知不对称持续（oracle $\rho{=}0.131{\pm}0.002$）。MHPR 在此规模无法分离 oracle（$0.187{\pm}0.157$）与 NoUnlearn（$0.152{\pm}0.096$）——其方差随类数增大——RII 同样无法分离它们（oracle $0.131{\pm}0.002$ 对比 NoUnlearn $0.127{\pm}0.026$）。此规模下稳健的信号是 RII 将 NegGrad 隐藏失败（$0.189{\pm}0.016$）与 NoUnlearn 分离。报告六种方法；SISA 因 20 类的分片训练在此规模不成比例而省略。

\begin{table}[t]
\centering
\small
\caption{CIFAR-100 上的类级遗忘基准（20 个训练类，遗忘类~3，$K{=}3$ 留出 $\{20,21,22\}$）；三个种子的均值 $\pm$ 标准差。}\label{tab:cifar100}
\setlength{\tabcolsep}{2pt}
\begin{tabular}{l c c c c c}
\toprule
方法 & Ret.(\%) & For.(\%) & $\rho$ & MHPR & MIA-loss \\
\midrule
NoUnlearn & $86.6{\pm}2.0$ & $74.3{\pm}13.1$ & $0.127{\pm}0.026$ & $0.152{\pm}0.096$ & $0.276{\pm}0.093$ \\
Retrain (oracle) & $85.3{\pm}2.0$ & 0.0 & $0.131{\pm}0.002$ & $0.187{\pm}0.157$ & 0.000 \\
NegGrad & $81.9{\pm}2.7$ & $21.7{\pm}4.1$ & $\mathbf{0.189{\pm}0.016}$ & $0.306{\pm}0.041$ & $0.096{\pm}0.019$ \\
FineTune & $92.5{\pm}0.4$ & $39.5{\pm}5.8$ & $0.178{\pm}0.004$ & $0.235{\pm}0.055$ & $0.077{\pm}0.015$ \\
KED & $93.1{\pm}0.5$ & $39.9{\pm}5.9$ & $0.174{\pm}0.005$ & $0.278{\pm}0.045$ & $0.068{\pm}0.015$ \\
BadTeacher & $93.1{\pm}0.6$ & $33.6{\pm}9.2$ & $0.160{\pm}0.003$ & $0.204{\pm}0.012$ & $0.063{\pm}0.020$ \\
\bottomrule
\end{tabular}
\end{table}

上述基准是视觉的。为检验跨模态迁移，我们在文本上重复该协议。在 AG News（四类，遗忘集 $=$ World 新闻）上，微调两轮 epoch 的 DistilBERT 分类器~\cite{sanh2020distilbert} 显示与视觉相同的排序：RII 将未触动模型（$\rho{=}0.267$）与重训练 oracle（$\rho{=}0.090$）区分开，$2.96{\times}$ 差距与 CIFAR-10 结果一致。这些文本结果是单次运行的 pilot 观察，作为迁移证据而非完整的跨模态基准报告，因为我们未在多个种子上重复。梯度上升在此行为不同：它彻底摧毁模型（保留准确率降到 $37.5\%$，$\rho{\approx}0$），是显性而非隐藏失败。因此低 RII 是擦除的必要而非充分证据：应与保留准确率一起解读，才能区分真移除与模型崩溃，这与第~\ref{sec:mia}~节的两轴观点一致。同样的情况出现在小语言模型（Qwen2.5-0.5B~\cite{qwen2024qwen2}）的作者级 TOFU~\cite{maini2024tofu} 协议下：暴露后在保留集上微调使遗忘答案保持不变（准确率不变），在仅解码器模型中复现第~\ref{sec:rebound}~节的反弹，而梯度上升再次摧毁模型（$\rho{\to}0$，损失增长一个数量级）。这个小模型 pilot 也暴露出一个真实的边界：在 0.5B 规模下，词表级输出通道无法将 oracle 与未触动模型区分开（$\rho{=}0.080$ 对比 $0.083$），因为两个答案集诱导出统计上相似的文本分布。第~\ref{sec:llm}~节表明在 7B 规模下相同词表通道变得可判别：证书的判别力在于输出通道的\emph{类语义}，而模型容量决定原始词表通道是否足够（参见注记~\ref{rem:channel}）；语言模型审计应优先使用类级通道（例如作者分类头）。

\subsection{动态范围与梯度扫描}

\begin{table}[t]
\centering
\small
\caption{RII 动态范围（CIFAR-10，cat 类）。}\label{tab:dynamic_range}
\setlength{\tabcolsep}{2pt}
\begin{tabular}{l c c}
\toprule
场景 & $\rho$ & $\sigma_2/\sigma_1$ \\
\midrule
随机 5\%（NoUnlearn） & $1.20\times10^{-3}$ & 0.035 \\
cat 类（基线）   & $\mathbf{1.83\times10^{-1}}$ & 0.47 \\
\bottomrule
\end{tabular}
\end{table}

表~\ref{tab:dynamic_range} 显示类遗忘的 $\rho$ 比随机遗忘\textbf{大 152$\times$}，表明 RII 并不退化。RII 在第一步的短暂上升之后随梯度上升步数下降（表~\ref{tab:gradient_sweep}），与式~\eqref{eq:finetune_bound}（定理~\ref{thm:finetune}）的 $O(\lambda^2)$ 界一致，其中 $\lambda$ 为上升步长。这些受控扫描（表~\ref{tab:dynamic_range}、表~\ref{tab:gradient_sweep}，以及表~\ref{tab:rebound} 的反弹研究）使用与第~\ref{sec:benchmark}~节主基准分开训练的基线，因此 NoUnlearn 参照值在此不同（$\rho{=}0.183$ 对比表~\ref{tab:benchmark} 的 $0.190{\pm}0.036$）；我们在这些研究中比较的是表内相对模式。

\begin{table}[t]
\centering
\small
\caption{RII 随梯度上升步数的变化（CIFAR-10，cat 类）。}\label{tab:gradient_sweep}
\setlength{\tabcolsep}{2pt}
\begin{tabular}{c c c c}
\toprule
步数 & $\rho$ & Acc.(\%) & $\sigma_2/\sigma_1$ \\
\midrule
0 & $1.83\times10^{-1}$ & 76.2 & 0.47 \\
1 & $1.95\times10^{-1}$ & 71.9 & 0.48 \\
3 & $1.76\times10^{-1}$ & 71.2 & 0.46 \\
10 & $1.05\times10^{-1}$ & 69.2 & 0.34 \\
30 & $1.25\times10^{-2}$ & 58.4 & 0.11 \\
100 & $2.04\times10^{-3}$ & 25.4 & 0.05 \\
\bottomrule
\end{tabular}
\end{table}

\subsection{MHPR 评估}

\begin{table}[t]
\centering
\small
\caption{MHPR 对比 LOCO（MNIST 类遗忘，在 8 类上训练，遗忘类~5）。}\label{tab:mhpr}
\setlength{\tabcolsep}{2pt}
\begin{tabular}{c c c c}
\toprule
$K$ & $\rho_{\mathrm{std}}$ & $\rho_{\mathrm{LOCO}}$ & $\rho_H$ \\
\midrule
1 & 0.145 & 0.379 & 0.942 \\
2 & 0.141 & 0.348 & 0.789 \\
\textbf{3} & 0.136 & 0.310 & \textbf{0.046} \\
4 & 0.142 & 0.167 & \textbf{0.0008} \\
\bottomrule
\end{tabular}
\end{table}

表~\ref{tab:mhpr} 比较 MHPR 与 LOCO。MHPR 将遗忘类均值投影到 $K$ 个留出未见类的张成子空间上（第~\ref{sec:mhpr}~节）；在 MNIST 类遗忘上，$K\ge3$ 解决 LOCO 失败：$K{=}1$ 给出 $\rho_H{=}0.942$（比 LOCO 更差，因为单个参照错失未见类流形），而 $K{=}3$ 将 $\rho_H$ 降到 $0.046$，比 $\rho_{\mathrm{std}}{=}0.136$ 与 $\rho_{\mathrm{LOCO}}{=}0.310$ 低一个数量级。$K{=}4$ 达到 $0.0008$（近机器零）但消耗训练数据（准确率 $79.5\%{\to}48.7\%$）。我们推荐 $K{=}3$。

跨数据集的自适应温度缩放将 MHPR 扩展到留出子空间可能坍缩的低精度区间（Fashion-MNIST 上一种种子依赖的失败，第~\ref{sec:benchmark}~节）。在 Fashion-MNIST（$88.7\%$ 准确率）的坍缩运行上，$T{=}50$ 将 $\rho_H$ 从 $\sim1$ 降到 $0.021$；在 CIFAR-10（$75\%$）上，$T{=}5$ 给出 $\rho_H{=}0.025$。这些优于类级 RII 基线（Fashion-MNIST 与 CIFAR-10 分别为 $0.166{\pm}0.004$ 与 $0.190{\pm}0.036$，见表~\ref{tab:cross}）。

\paragraph{理想遗忘构造。}为进一步验证 MHPR，我们构造一个\emph{精确}位于 $K{=}5$ 个留出 Fashion-MNIST 类均值张成子空间内的合成遗忘均值 $\boldsymbol{\mu}_f$。所得 $\rho_H$ 为机器零（$<10^{-10}$）。注入各向同性噪声 $\boldsymbol{\varepsilon}\sim\mathcal{N}(0,\varepsilon^2\mathbf{I})$ 使 $\rho_H$ 单调上升：$\varepsilon{=}0.01$ 给出 $\rho_H{=}0.005$，$\varepsilon{=}0.10$ 给出 $\rho_H{=}0.237$。这证实 $\rho_H{=}0$ 当且仅当 $\boldsymbol{\mu}_f$ 位于留出子空间内，直接验证命题~\ref{prop:mhpr_monotone}。

\subsection{微调反弹与敏感性}
\label{sec:rebound}

第~\ref{sec:benchmark}~节将梯度上升排为最弱的类级遗忘方法。一个实际问题：在实践常用的"先上升后微调"配方下，该排序是否仍然成立。

\paragraph{微调反弹。}当在 $D_f$ 上的梯度上升之后接着在 $D_r$ 上微调时，$\rho$ 从 $0.105$（仅上升）升到 $0.221$（上升 $+$ 微调），甚至超过基线 $0.183$（表~\ref{tab:rebound}）。定理~\ref{thm:unified} 解释这一点：上升降低覆盖泄漏 $\eta^2$，但微调重塑决策边界并使遗忘表示重新错位，抬升 $\sigma_2$——这是基于准确率的指标会错过的失败。

\begin{table}[t]
\centering
\small
\caption{微调反弹现象（CIFAR-10，cat 类）。}\label{tab:rebound}
\setlength{\tabcolsep}{2pt}
\begin{tabular}{l c c c}
\toprule
操作 & $\rho$ & Acc.(\%) & $\sigma_2/\sigma_1$ \\
\midrule
基线（不遗忘） & 0.183 & 76.2 & 0.47 \\
上升 10 步 & 0.105 & 69.2 & 0.34 \\
上升 10 + 微调 & \textbf{0.221} & 75.2 & 0.52 \\
\bottomrule
\end{tabular}
\end{table}

\paragraph{跨数据集与敏感性。}对 $C{=}10$--$100$（MNIST、CIFAR-10/100），$\rho\le2\times10^{-3}$，证实秩一性质是结构性的。在刻意过拟合下（表~\ref{tab:sensitivity}，第~\ref{app:extra}~节），$\rho$ 增加十倍（MNIST 60K$\to$5K，50 epoch：$9.8\times10^{-4}\to1.24\times10^{-2}$）。受控 $\alpha$-扰动证实正交轴：$\alpha{=}0.20$ 使 RII 提高 $9\times$ 而 MIA 仅上升 $6.8$ 个百分点。CIFAR-100 上的 ResNet-18 保持 $\rho{=}1.41\times10^{-3}$（表~\ref{tab:resnet18}），证实架构独立性。

\paragraph{逐样本 MHPR 差距。}命题~\ref{prop:ps_mhpr} 预测 $\bar{\rho}_H \ge \|\boldsymbol{\pi}_f\|_2^2\,\rho_{H,\mathcal{S}}$（Jensen 差距）。在我们的 MNIST 实验中，基线时 $\rho_{H,\mathcal{S}}{\approx}0.99999$ 且 $\bar{\rho}_H{\approx}0.99932$，差距约 $0.00066$，在 100 个梯度步后缩小到 $0.00023$；在更多异质数据上预期差距更大，使逐样本 MHPR 更保守。

\subsection{大规模验证与敏感性}
\label{app:extra}

\begin{table}[t]
\centering
\small
\caption{CIFAR-100 上的 ResNet-18。}\label{tab:resnet18}
\setlength{\tabcolsep}{3pt}
\begin{tabular}{l c c c c}
\toprule
模型 & $|D_f|$ & Acc.(\%) & $\rho$ & $\sigma_2/\sigma_1$ \\
\midrule
SmallCNN & 2,500 & 20.8 & $7.10{\times}10^{-4}$ & 0.027 \\
ResNet-18 & 2,500 & 29.7 & $1.41{\times}10^{-3}$ & 0.038 \\
\bottomrule
\end{tabular}
\end{table}

\begin{table}[t]
\centering
\small
\caption{RII 对过拟合的敏感性。}\label{tab:sensitivity}
\setlength{\tabcolsep}{2pt}
\begin{tabular}{l l c c c c}
\toprule
数据集 & 模型 & $|D|$ & Ep. & Acc.(\%) & $\rho$ \\
\midrule
MNIST   & SimpleMLP & 60K & 10  & 97.45 & $9.8{\times}10^{-4}$ \\
MNIST   & DeepMLP & 5K & 50 & 95.50 & $\mathbf{1.24{\times}10^{-2}}$ \\
CIFAR-10 & SmallCNN & 50K & 10  & 76.34 & $1.5{\times}10^{-3}$ \\
CIFAR-10 & SmallCNN & 5K & 50 & 61.63 & $\mathbf{1.31{\times}10^{-2}}$ \\
\bottomrule
\end{tabular}
\end{table}

在刻意过拟合下（表~\ref{tab:sensitivity}），$\rho$ 增加十倍，证实对记忆化的敏感性。在 10 次独立 MNIST 训练运行（5\% 随机遗忘）中，$\rho$ 均值 $9.72{\times}10^{-4}$、标准差 $2.1{\times}10^{-5}$（CV $2.2\%$），MIA 均值 $90.4\%$、标准差 $0.3\%$；两者都稳定但正交，如定理~\ref{thm:unified} 所预测。

\subsection{与现有指标的对比}

\begin{table}[t]
\centering
\caption{遗忘评估指标对比。}\label{tab:metric_compare}
\setlength{\tabcolsep}{3pt}
\begin{tabular}{l c c c}
\toprule
指标 & 信息论 & 免调参 & 复杂度 \\
\midrule
MIA~\cite{shokri2017membership} & 否 & 否 & 中 \\
DP~\cite{guo2020certified} & 是 & 否 & 高 \\
Backdoor~\cite{sommer2022towards} & 否 & 否 & 高 \\
表示级~\cite{cosma2026ruler} & 否 & 否 & 高 \\
审计~\cite{ye2026auditing} & 否 & 否 & 高 \\
\textbf{RII} & \textbf{是} & \textbf{是} & $O(CN)$ \\
\textbf{MHPR} & \textbf{是} & \textbf{是} & $O(KCN)$ \\
\bottomrule
\end{tabular}
\end{table}

表~\ref{tab:metric_compare} 将 RII 与 MHPR 和现有指标对比。RII 是表中唯一同时满足谱奠基、免超参计算与 $O(CN)$ 复杂度的条目。

% =============================================================================
% zh_part3.tex — §6–§9（逐句对照中文版）
% =============================================================================

% =============================================================================
\section{调和 MIA 与 RII：两个正交泄漏轴}
\label{sec:mia}

在随机遗忘基线中（表~\ref{tab:main}），训练/测试成员推断准确率保持在约 90\%，而 $\rho \approx 10^{-3}$。这两者并不矛盾：此处的 MIA 度量一般训练成员性——训练样本与测试样本之间的置信度差距——而 RII 度量遗忘与保留输出是否可区分。对随机遗忘子集，模型并未改变，因此遗忘与保留输出仍不可区分（$\rho\approx10^{-3}$），尽管训练/测试成员信号与任何已训练模型一样强（90\%）。定理~\ref{thm:unified} 将这两者形式化为两条正交的泄漏轴。

\begin{itemize}
\item \textbf{分布位移（MIA 轴，$I_{\mathrm{dir}}$）：}任何训练样本往往比测试样本有更高置信度。MIA 利用这一差距，区分 $P(\ell\mid\text{train})$ 与 $P(\ell\mid\text{test})$。
\item \textbf{样本特定泄漏（RII 轴，$I_{\mathrm{cov}}$）：}RII 度量 $P(Y\mid D_f,\theta') \neq P(Y\mid D_r,\theta')$ 是否成立，即 $D_f$ 是否留下不同的输出签名。
\end{itemize}

Retrain 与 NoUnlearn 可表现出几乎相同的 RII，同时两者 MIA 准确率都高，因为 MIA 检测的是训练成员性的一般性，而非某个具体遗忘请求是否被履行。\textbf{RII 证明输出级遗忘/保留不可区分性}——这正是擦除合规的相关保证。

\begin{table}[t]
\centering
\caption{展示正交泄漏轴的受控扰动实验（MNIST，5\% 遗忘）。}\label{tab:alpha}
\setlength{\tabcolsep}{3pt}
\begin{tabular}{c c c c}
\toprule
$\alpha$ & $\rho$ (RII) & MIA (\%) & 解释 \\
\midrule
0.00 & $9.76{\times}10^{-4}$ & 90.4 & 基线 \\
0.20 & $8.91{\times}10^{-3}$ & 97.2 & 强直接泄漏 \\
\bottomrule
\end{tabular}
\end{table}

表~\ref{tab:alpha} 证实定理~\ref{thm:unified} 预测的不对称响应：随着 $\alpha$ 增加，RII 上升 $9\times$ 而 MIA 增加 6.8 个百分点。这种不对称直接反映泄漏的独立轴。相同正交性出现在真实遗忘方法下（第~\ref{sec:benchmark}~节）：RII 与 MIA-loss 相关 $r{=}0.57$（中等，基于七个方法级均值），且顶端排序反转：NegGrad 在 RII 下最差、在 MIA-loss 下第二差。

\paragraph{超越输出。}
输出层保证对白盒安全是必要而非充分的。利用链式法则 $I(D_f;Y,\theta') = I(D_f;\theta') + I(D_f;Y\mid\theta')$，总泄漏分解为 $I(D_f;\theta') = I(D_f;Y) + I(D_f;\theta'\mid Y) - I(D_f;Y\mid\theta')$。RII 证明条件项 $I(D_f;Y\mid\theta')$（给定模型的输出泄漏）。项 $I(D_f;\theta'\mid Y)$（输出之外的内部表示中的信息）可以通过审计内部特征（例如倒数第二层激活上的 MMD）来界定。两者共同构成完整的两层证书；完整刻画留待未来工作。经验上（第~\ref{sec:benchmark}~节），倒数第二层特征上的残差探针甚至对重训练 oracle 也能以 AUC ${\approx}0.96$ 分离遗忘与留出样本，证实输出级证书必须与表示级审计配对。

\paragraph{与现有指标对比。}
在所比较的指标中，RII 与 MHPR 是唯一同时满足谱奠基、免超参计算与尺度不变性的。与 MIA（混淆训练/测试位移与遗忘/保留泄漏，$O(CN)$ 但依赖调参）、DP 认证（侵入式训练）、或基于后门的审计（依赖攻击）不同，RII 直接量化被删数据的影响是否已从模型的输出分布中移除（表~\ref{tab:metric_compare}）。

% =============================================================================
\section{相关工作}
\label{sec:related}

\textbf{机器学习遗忘。}该领域涵盖精确重训练~\cite{cao2015towards}、SISA~\cite{bourtoule2021machine}、基于梯度的反转~\cite{golatkar2020eternal}、基于删除的方法~\cite{ginart2019making,sekhari2021remember}，以及影响函数~\cite{guo2020certified}。我们的贡献对任何具体算法都是正交的：一个度量输出空间不可逆性的算法无关\emph{评估指标}。

\textbf{遗忘评估。}成员推理攻击~\cite{shokri2017membership,carlini2022membership} 是事实标准，但 MIA 检测训练成员性的一般性而非遗忘集特定的泄漏；差分隐私界~\cite{guo2020certified} 给出最坏情况保证但要求侵入式训练；后门审计依赖攻击。RII 则是一个免调参的输出级不可区分性谱证书。PSSD（附录~\ref{app:pssd}）共享这一哲学：$\delta(x)$ 得分是一个证书而非攻击——如果所有遗忘样本的 $\delta(x)$ 都低，那么任何基于阈值的 MIA 都无法将它们与保留样本分离（定理~\ref{thm:pssd_mia}），这是比经验攻击准确率更强的操作保证。

\textbf{超越攻击的验证。}近期工作表明模型能通过输出级检查、同时在中间表示中保留残余签名~\cite{cosma2026ruler,yong2026erased}；表示级验证~\cite{cosma2026ruler} 与通用审计框架~\cite{ye2026auditing} 在不使用影子模型或重训练基线的情况下解决此问题，而~\cite{zhang2024fragile} 表明验证可被不诚实提供者绕过，~\cite{zheng2026bias} 分析类级遗忘中的分类头偏差；我们请读者参阅~\cite{xue2026verification} 获取分类体系。RII 是这些工作的补充：一个纯输出级、免调参的谱证书，不需要影子模型或训练时修改。

\textbf{信息论安全。}我们的谱分析建立在 Shannon 的信道模型~\cite{shannon1949communication} 与现代信息论安全~\cite{bloch2021overview} 之上：Hilbert--Schmidt 分解给出我们 $2\times C$ 分析的连续类比，$\chi^2$-MI 界是标准的。奇异值衰减与残余输出泄漏之间的联系（定理~\ref{thm:unified}）将信息论隐私扩展到遗忘，其中``秘密''是遗忘集成员性而非固定属性。

% =============================================================================
\section{部署考量}
\label{sec:deployment}

RII 需要对 $D_f \cup D_r$ 做一次前向传播，花费 $O(C(N_f+N_r))$ 时间与 $O(C)$ 内存，相比模型训练可忽略。对大规模遗忘集，命题~\ref{prop:confidence} 提供子采样保证：从 $D_f$ 与 $D_r$ 各取
\[
N_{\min} \ge \frac{16C\nu^2}{\sigma_1^2\varepsilon^2}\log\frac{4C}{\delta}
\]
个样本，则估计满足 $|\hat{\rho}-\rho^*| \le \varepsilon$ 且置信度 $\ge 1-\delta$。取 $\sigma_1\approx0.4$、$\nu\approx1$、$C=10$，以绝对精度 $\varepsilon=10^{-2}$ 解析 $\rho$ 在最坏情况需 $N_{\min}\approx7\times10^7$；该界是保守的（对 $C$ 个坐标的 union bound 与均值常数），实际中每类约 $10^4$ 样本（如我们的 MNIST 规模审计）即可解析相关谱区间。MHPR 需要对留出类额外做 $K$ 次前向传播。

\paragraph{推荐审计协议。}
(1) 从 $D_f$ 与 $D_r$（以及 MHPR 的每个留出类）均匀抽取 $N_{\min}$ 个样本。(2) 计算 $\hat{\rho}$（或 $\hat{\rho}_H$），若模型准确率 $\lesssim 90\%$ 则应用温度缩放。(3) 报告 $[\hat{\rho}-\varepsilon,\;\hat{\rho}+\varepsilon]$ 作为水平 $1-\delta$ 的保守高概率区间。(4) 对两轴评估，报告 RII 旁的 MIA AUC，通过 $I(D_f;\theta') \lesssim I_{\mathrm{RII}} + I_{\mathrm{MIA}}$ 界定总泄漏。

\paragraph{估计直接泄漏参数 $\alpha$。}
统一分解将泄漏分为 $I_{\mathrm{cov}}\propto\eta^2$（由 RII 测量）与 $I_{\mathrm{dir}}\propto\alpha^2$。$\eta$ 可直接从通道矩阵计算，而 $\alpha$ 在黑盒审计中不可观测；MIA 相对随机猜测的优势提供一个下界，因为 $\mathrm{AUC}>0.5$ 蕴含 $\alpha>0$，且差距 $\mathrm{AUC}-0.5$ 随 $|\alpha|$ 单调缩放。因此两轴审计（RII 测覆盖泄漏，MIA AUC 测直接泄漏）通过 $I(D_f;\theta') \lesssim I_{\mathrm{RII}} + I_{\mathrm{MIA}}$ 界定总输出泄漏，无需直接观测 $\alpha$：RII 证明遗忘集的输出签名已被擦除，而 MIA AUC 界定黑盒攻击者可利用的残余分布位移。

% =============================================================================
\section{结论}
\label{sec:conclusion}

类级遗忘验证处于两种失败模式之间：基于准确率的指标在高遗忘强度下饱和并失去排序力，而 MIA 度量一个相关但不同的量——一般训练成员性——而非遗忘集特定的输出泄漏。RII 填补这一空白：它是一个免调参、黑盒的输出级不可区分性证书，恰好在准确率饱和时仍具信息量。经验上，它在大多数设置中跟踪已知遗忘强度并标记效用驱动的崩溃，并在 LLaMA-2-7B 的 TOFU 基准上在生成准确率饱和于 100\% 的检查点中保持该排序。在 MUSE 风格长文本书籍遗忘上（第~\ref{sec:muse}~节），梯度上升同样使遗忘输出分布保持不变，遗忘困惑度将重训练排序为最强擦除，而原始词表通道反映的是两本书的固有差异而非遗忘状态——这是强调通道语义作用的一个清晰边界情形。实验也划定了谱证书能做什么、不能做什么。RII 在梯度上升下\emph{上升}（即便准确率与 MIA 都改善，且随类数增加而增强），在上升-后-微调后反弹，并在模型退化时趋于零，因此必须始终与效用指标一起解读。在理论侧，$\rho=0$ 等价于输出通道独立，微调界为 $O(\lambda^2)$，统一分解将泄漏分为直接（MIA）与覆盖（RII）两轴。MHPR 扩展及其温度缩放处理已知/未知类不对称，PSSD（附录~\ref{app:pssd}）将审计扩展到逐样本状态差异。

\textbf{局限性。}该证书按构造是输出级的：它不界定内部表示中的泄漏（探针实验表明即使对重训练 oracle 这种泄漏也可能存在），也不证明参数级隐私。类级基准基于三个种子（以均值 $\pm$ 标准差报告），它们支持指标的稳定性与定性排序，但在 $n{=}3$ 下并不单独确立相近对的统计显著性。随机基线（10 种子）与 7B 研究（3 次子采样）上的种子稳定性进一步佐证指标的稳定性。AG News、0.5B 与 MUSE 风格书籍结果是单次运行的 pilot 观察，作为迁移证据而非完整基准报告（书籍研究还受基座模型对公有领域文本的先前暴露限制，这压缩了 teacher 的可学习余量）。MI 界要求 $\pi_{\min}>0$，统一分解在高斯信号模型下精确成立、一般情形一阶成立。

\textbf{实用指南。}审计中不应单独用 RII 来宣布遗忘成功。由于该证书是输出级的，且既可能因模型退化也可能因擦除而趋于零，我们建议同时报告：(i) 效用检查（保留准确率，以及可行时的遗忘准确率），以区分真擦除与崩溃；(ii) 一个 MIA 得分（遗忘/保留损失 AUC，或适当的训练/测试成员准确率），用于直接泄漏轴；(iii) 内部表示探针（倒数第二层特征上的 MMD 或残差线性探针），以捕捉表示级签名；(iv) 在留出参照类可用时使用 MHPR。第~\ref{sec:mia}~节形式化的正是这种两层、两轴审计。

\textbf{未来工作。}方向包括正式的 DP--RII 联系；将作者级 TOFU 研究扩展到完整基准套件（遗忘 90/95/99、Truth Ratio 与 KLR~\cite{maini2024tofu}），并将 MUSE 风格书籍研究（第~\ref{sec:muse}~节）扩展到完整 MUSE~\cite{shi2024muse} 套件（版权书籍、新闻）；语言模型的类级输出通道（例如作者或书籍分类头）；PSSD 的多层状态轨迹扩展（附录~\ref{app:pssd}）；将 RII 与内部表示审计结合以实现完全白盒认证；ImageNet 规模验证；以及在缺乏语义匹配的留出类时，用合成或生成式参考分布构造 MHPR 参考子空间。

% =============================================================================
% zh_part4.tex — 附录 A/B/C + 参考文献（逐句对照中文版）
% =============================================================================

\appendix

% =============================================================================
\section{PSSD：逐样本状态差异扩展}
\label{app:pssd}

RII 与 MHPR 证明通道的输出级不可区分性，但无法检测异常的\emph{个别}样本（注记~\ref{rem:channel}）。本附录发展 PSSD（逐样本状态差异，Per-Sample State Disparity），一种"先度量后平均"的改进，为每个样本计算差异得分。论文的核心 RII/MHPR 贡献不依赖此扩展。

\subsection{定义}

设 $Q:\mathbb{R}^C\to\{1,\dots,m\}$ 为第~\ref{sec:logit_bridge}~节的 $k$-means 量化器。\textbf{保留参照分布}为
\begin{equation}
\pi_r(s) = \frac{1}{N_r}\sum_{x\in D_r}\mathbf{1}\{Q(z(x))=s\}, \qquad \boldsymbol{\pi}_r\in\Delta^{m-1}.
\end{equation}

\begin{definition}[个体状态差异（ISD）]
\label{def:isd}
对状态为 $s(x)=Q(z(x))$ 的 $x$，$\delta(x)\coloneqq 1-\pi_r(s(x))\in[0,1-1/m]$，度量样本的 logit 状态相对保留分布的``令人惊讶''程度。
\end{definition}

\begin{definition}[聚合 PSSD 指标]
\label{def:pssd_agg}
\begin{align}
\textbf{MPSD:}\quad &\Delta_f \coloneqq \frac{1}{N_f}\sum_{x\in D_f}\delta(x), &
\textbf{超出量:}\quad &\Psi \coloneqq \Delta_f - \Delta_r,\\
\textbf{Top-$k$:}\quad &\Delta_{(k)} \coloneqq \frac{1}{k}\sum_{i=1}^k\delta_{[i]},
\end{align}
其中 $\delta_{[1]}\ge\cdots$ 为次序统计量。$\Psi$ 量化遗忘样本超出基线水平的\emph{额外}惊讶程度，而 $\Delta_{(k)}$ 刻画最坏情况的个体泄漏。
\end{definition}

\subsection{理论保证}

\begin{proposition}[PSSD--RII 联系]
\label{prop:pssd_rii}
$\Delta_f = 1-\langle\boldsymbol{\pi}_f,\boldsymbol{\pi}_r\rangle$。当 $\|\boldsymbol{\pi}_f\|_2=\|\boldsymbol{\pi}_r\|_2=s$ 时，状态空间 RII 满足 $\rho_{\mathcal{S}}=\frac{s-\langle\boldsymbol{\pi}_f,\boldsymbol{\pi}_r\rangle}{2s}=\frac{1}{2}\bigl(1-\frac{\langle\boldsymbol{\pi}_f,\boldsymbol{\pi}_r\rangle}{\|\boldsymbol{\pi}_f\|_2^2}\bigr)$。差距 $|\Delta_f-\rho_{\mathcal{S}}|\le\|\boldsymbol{\pi}_f-\boldsymbol{\pi}_r\|_1(1-1/m)+O(\|\boldsymbol{\pi}_f-\boldsymbol{\pi}_r\|_2^2)$。因此当 $\boldsymbol{\pi}_f\to\boldsymbol{\pi}_r$ 时 $\Delta_f$ 与 $\rho_{\mathcal{S}}$ 渐近等价，且 $\rho_{\mathcal{S}}=0$ 当且仅当 $\Delta_f = 1-\|\boldsymbol{\pi}_r\|_2^2$。
\end{proposition}

\begin{proof}[证明梗概]
$\Delta_f=\frac1{N_f}\sum_i(1-\pi_r(s_i))=1-\langle\boldsymbol{\pi}_f,\boldsymbol{\pi}_r\rangle$。差距项通过围绕 $\boldsymbol{\pi}_f=\boldsymbol{\pi}_r$ 展开 $\rho_{\mathcal{S}}$ 并界定 $|\langle\boldsymbol{\pi}_f,\boldsymbol{\pi}_r\rangle-\|\boldsymbol{\pi}_r\|_2^2|\le\|\boldsymbol{\pi}_f-\boldsymbol{\pi}_r\|_1(1-1/m)$ 得到。
\end{proof}

\begin{proposition}[集中界]
\label{prop:pssd_conc}
由于 $\delta(x)\in[0,1]$，Hoeffding 给出 $|\hat{\Delta}_f-\Delta_f|\le\sqrt{\log(2/\zeta)/(2N_f)}$ 与 $|\hat{\Psi}-\Psi|\le\sqrt{2\log(4/\zeta)/N_{\min}}$，置信度 $\ge1-\zeta$。与命题~\ref{prop:confidence} 不同，无需次高斯假设。
\end{proposition}

\begin{proposition}[逐样本 MHPR Jensen 差距]
\label{prop:ps_mhpr}
定义逐样本状态残差 $\rho_H(x)\coloneqq\min_{\mathbf{w}}\|\mathbf{e}_{s(x)}-\sum_k w_k\boldsymbol{\pi}_{h_k}\|_2^2$ 与 $\bar{\rho}_H\coloneqq\frac1{N_f}\sum_x\rho_H(x)$。则 $\bar{\rho}_H \ge \|\boldsymbol{\pi}_f\|_2^2\,\rho_{H,\mathcal{S}}$；特别地 $\bar{\rho}_H\ge0$，当且仅当遗忘类是状态齐次时取等。该差距等于遗忘类投影后类内状态方差的平均值。
\end{proposition}

\begin{proof}[证明梗概]
对固定子空间 $\mathcal{S}_H$，残差 $\|\mathbf{v}-\mathbf{P}_{\mathcal{S}_H}\mathbf{v}\|_2^2$ 是关于 $\mathbf{v}$ 的凸二次型。Jensen 给出 $\frac1{N_f}\sum_x\rho_H(x)\ge\rho_H(\frac1{N_f}\sum_x\mathbf{e}_{s(x)})=\rho_H(\boldsymbol{\pi}_f)$，且 $\rho_H(\boldsymbol{\pi}_f)=\min_{\mathbf{w}}\|\boldsymbol{\pi}_f-\sum_k w_k\boldsymbol{\pi}_{h_k}\|_2^2=\|\boldsymbol{\pi}_f\|_2^2\,\rho_{H,\mathcal{S}}$。精确差距为 $\frac1{N_f}\sum_x\|(\mathbf{I}-\mathbf{P}_{\mathcal{S}_H})(\mathbf{e}_{s(x)}-\boldsymbol{\pi}_f)\|_2^2\ge0$，仅当遗忘类是状态齐次（所有样本映射到同一状态）时消失。
\end{proof}

\begin{theorem}[基于 PSSD 的 MIA]
\label{thm:pssd_mia}
分类器 $T_\tau(x)=\mathbf{1}\{\delta(x)\ge\tau\}$ 满足：(i) $\Psi>0\Rightarrow\mathrm{AUC}>0.5$；(ii) $\mathrm{AUC}\ge1-\frac12\mathrm{H}^2(P_{\delta|f},P_{\delta|r})$（Hellinger 界）；(iii) $|\widehat{\mathrm{AUC}}-\mathrm{AUC}|\le\sqrt{\log(2/\zeta)/(2(N_f+N_r))}$。
\end{theorem}

\paragraph{状态轨迹（扩展）。}
对 $L$ 层模型，每个样本有状态轨迹 $\mathbf{s}(x)=(s_1,\dots,s_L)$；将其建模为一阶马尔可夫链可将 PSSD 扩展到单层 logits 之外，留待未来工作。完整证明见附录~\ref{sec:appendix}。

\subsection{实验结果}

我们在类级遗忘上评估 PSSD，使用 $k$-means 量化（$m{=}20$ 状态，$K{=}3$ 留出）。表~\ref{tab:pssd_results} 报告 5 个梯度上升步后 MPSD $\Delta_f$、超出量 $\Psi$，以及 $\delta$ 基与置信基两种分类器的 MIA AUC。

\begin{table}[t]
\centering
\small
\caption{PSSD 指标对比 MIA（$m{=}20$，$K{=}3$，5 步上升）。}\label{tab:pssd_results}
\setlength{\tabcolsep}{2pt}
\begin{tabular}{l c c c c}
\toprule
数据集 & $\Delta_f$ & $\Psi$ & MIA($\delta$) & MIA(conf.) \\
\midrule
MNIST & 0.997 & 0.066 & \textbf{0.990} & 0.570 \\
Fashion-MNIST & 0.990 & 0.057 & \textbf{0.943} & 0.701 \\
CIFAR-10 & 0.978 & 0.001 & 0.511 & 0.807 \\
\bottomrule
\end{tabular}
\end{table}

$\delta$ 基 MIA 在 MLP 模型上优于置信基 MIA（表~\ref{tab:pssd_results}），利用 softmax 置信度随机时仍可判别的 logit 几何。对 CIFAR-10（CNN），优势显著收窄：$\mathrm{AUC}(\delta){=}0.511$ 对比 $\mathrm{AUC}(\mathrm{conf.}){=}0.807$，因为 CNN 的 softmax 置信度本身提供强 MIA 信号，而超出差异 $\Psi{=}0.001$ 近零。

\paragraph{对 $m$ 与量化的敏感性。}在 MNIST 上跨 $m{=}20,50,100$，MPSD 稳定（$\Delta_f{=}0.997{\pm}0.002$）且 $\delta$-MIA AUC 保持 $>0.97$，而 $\Psi$ 从 $0.066$ 降到 $0.014$，$\rho_{\mathcal{S}}$ 从 $0.176$ 降到 $0.131$，随码本细化收敛到 $\rho_{\mathrm{sm}}{=}0.109$（定理~\ref{thm:state_bridge}）。我们推荐默认 $m{=}2C{=}20$。

% =============================================================================
\section{状态空间桥与算子理论推广}
\label{app:bridge}

\subsection{量化与谱桥}

\begin{lemma}[量化误差]
\label{lem:quantization}
对 $N$ 个独立同分布、$\|z\|_2\le R$ 的 logit 向量，带 $m$ 个中心的 $k$-means 给出量化误差 $\mathbb{E}[\|z-\tilde{z}\|_2] \le \tilde{O}(R\sqrt{Cm/N}) + O(R\sqrt{\log(1/\zeta)/N})$，置信度 $\ge1-\zeta$，其中 $\tilde{z}$ 为码本重构。证明见附录~\ref{sec:appendix}。
\end{lemma}

\begin{theorem}[谱桥：状态 $\leftrightarrow$ softmax]
\label{thm:state_bridge}
设 $\mathbf{a}_s = \mathbb{E}[\mathbf{p}(x) \mid Q(z(x))=s]$ 为状态 $s$ 内的平均 softmax，$\mathbf{A} = [\mathbf{a}_1,\dots,\mathbf{a}_m]^\top$ 的奇异值为 $\tau_1 \ge \cdots \ge \tau_m > 0$。设 $\mathbf{B}\in\mathbb{R}^{2\times m}$ 的两行为遗忘/保留的状态条件概率，则 softmax 通道分解为 $\mathbf{M}_{\mathrm{sm}}=\mathbf{B}\mathbf{A}$。假设 softmax 为 $L$-Lipschitz 且量化误差 $\le \delta_q$（引理~\ref{lem:quantization}）。则
\begin{equation}
|\rho_{\mathcal{S}} - \rho_{\mathrm{sm}}| \le 2\Bigl(\frac{\tau_1}{\tau_m}-1\Bigr)\rho_{\mathcal{S}} + O\!\Bigl(\Bigl(\frac{\tau_1}{\tau_m}-1\Bigr)^2 + \frac{L\delta_q}{\|\boldsymbol{\mu}_f\|_2}\Bigr).
\end{equation}
若码本良态（$\tau_1/\tau_m \le 1+\kappa_0$）且 $\delta_q$ 小，则两个 RII 值等价到常数因子。证明利用分解 $\mathbf{M}_{\mathrm{sm}}=\mathbf{B}\mathbf{A}$ 并通过右逆论证界定奇异值（附录~\ref{sec:appendix}）。
\end{theorem}

\begin{theorem}[MHPR 对 softmax RII 的偏差]
\label{thm:softmax_deviation}
设 $\boldsymbol{\mu}_f^*$ 为遗忘类的总体平均 softmax，$\mathbf{H}^* \in \mathbb{R}^{K\times C}$ 为 $K$ 个留出类总体均值的矩阵，$r=\mathrm{rank}(\mathbf{H}^*)\le K$。假设理想条件 $\boldsymbol{\mu}_f^* \in \mathrm{rowspan}(\mathbf{H}^*)$。设 $\hat{\boldsymbol{\mu}}_f$、$\hat{\mathbf{H}}$ 为来自 $N_f$ 与 $N_{h_1},\dots,N_{h_K}$ 个样本的经验估计，$N_{\min} = \min(N_f,N_{h_1},\dots,N_{h_K})$，$\sigma_{\min}(\mathbf{H}^*)$ 为 $\mathbf{H}^*$ 的最小奇异值。则置信度 $\ge 1-\delta$ 时，
\begin{equation}
\rho_H \le O\!\left(\frac{r\,\nu^2 C\,\log((K+1)C/\delta)}{\sigma_{\min}^2(\mathbf{H}^*)\|\boldsymbol{\mu}_f^*\|_2^2\,N_{\min}}\right) + O(N_{\min}^{-3/2}).
\end{equation}
因此 $\rho_H = O(1/N_{\min})$，快于标准 RII 界的 $O(1/\sqrt{N})$ 阶。
\end{theorem}

\subsection{算子理论推广}

离散 $2\times C$ 分析通过 Hilbert--Schmidt 算子理论提升到连续输入/输出空间，为定理~\ref{thm:unified} 提供基础。在温和正则性下，条件密度有典范分解 $f_{Y|X}(y|x)=f_Y(y)+\sum_{k=1}^{\infty}\sigma_k\,\phi_k(y)\,\psi_k(x)$，其中 $\sigma_1\ge\sigma_2\ge\cdots\ge0$，$\{\phi_k\},\{\psi_k\}$ 分别是 $L^2(\mathcal{Y}),L^2(\mathcal{X})$ 按 $P_Y,P_X$ 加权的标准正交基。第一项是边际密度（$\sigma_1=1$——加权基的归一化，区别于离散通道矩阵 $\mathbf{M}$ 的未加权首奇异值，后者不必等于 1）。\textbf{偏离参数} $\eta^2\coloneqq\sum_{k\ge2}\sigma_k^2$ 量化所有秩 $\ge2$ 分量的总能量；在离散 $2\times C$ 情形，$\eta^2=\sigma_2^2$ 且 $\rho=\eta^2/(\sigma_1^2+\eta^2)$，因此 $\rho$ 关于 $\eta^2$ \textbf{单调}，且 $\mathbb{E}_X[\chi^2(P_{Y|X}\|P_Y)]=\eta^2$ 确立 $\eta^2$ 既是谱量又是散度论量。

% =============================================================================
\section{关键证明与补充结果}
\label{sec:appendix}

\subsection{定理~\ref{thm:perfect} 的证明（输出分布不可逆性）}
\label{app:perfect}
\begin{proof}
我们证明两个方向。

\noindent\textit{（$\Rightarrow$）}$\rho=0 \Rightarrow \sigma_2=0 \Rightarrow \mathbf{M}$ 秩为 1。由于两行都是和为 1 的概率向量，秩一 $2\times C$ 矩阵的行必成比例：$\boldsymbol{\mu}_f = \lambda \boldsymbol{\mu}_r$ 且 $\lambda>0$。和约束 $\sum_c (\boldsymbol{\mu}_f)_c = \sum_c (\boldsymbol{\mu}_r)_c = 1$ 迫使 $\lambda=1$，故 $\boldsymbol{\mu}_f = \boldsymbol{\mu}_r$。由全概率公式，$(\boldsymbol{\mu}_f)_c = \frac{1}{N_f}\sum_{x\in D_f}p_c(x) = P(Y{=}c\mid X{=}f)$ 且 $(\boldsymbol{\mu}_r)_c = P(Y{=}c\mid X{=}r)$，所以 $\boldsymbol{\mu}_f=\boldsymbol{\mu}_r$ 蕴含 $P(Y\mid X=1,\theta') = P(Y\mid X=0,\theta')$，即 $Y\perp X\mid\theta'$ 且 $I(X;Y\mid\theta')=0$。无需指数族或校准假设：$\mathbf{M}$ 的行按构造就是通道的条件输出分布（注记~\ref{rem:channel}）。

\noindent\textit{（$\Leftarrow$）}若 $I(X;Y\mid\theta')=0$，则 $P(Y\mid X{=}f)=P(Y\mid X{=}r)$，即 $\boldsymbol{\mu}_f = \boldsymbol{\mu}_r$，故 $\sigma_2=0$ 且 $\rho=0$。
\end{proof}

\subsection{命题~\ref{prop:confidence} 的证明（有限样本置信）}
\label{app:confidence}
\begin{proof}
在次高斯假设~\cite{vershynin2018high} 下，每个坐标满足
\begin{equation}
\mathbb{P}\bigl(|(\hat{\boldsymbol{\mu}}_f)_c-(\boldsymbol{\mu}_f^*)_c|>t\bigr)\le 2\exp(-N_ft^2/2\nu^2),
\end{equation}
对 $C$ 个坐标做 union bound 给出 $\|\hat{\boldsymbol{\mu}}_f-\boldsymbol{\mu}_f^*\|_2\le\nu\sqrt{2C\log(4C/\delta)/N_f}$（置信度 $\ge1-\delta/2$），对 $\hat{\boldsymbol{\mu}}_r$ 同理；因此
\begin{equation}
\|\hat{\mathbf{M}}-\mathbf{M}^*\|_F\le2\nu\sqrt{C\log(4C/\delta)/N_{\min}}.
\end{equation}
由 Weyl 定理，$\|\hat{\boldsymbol{\sigma}}-\boldsymbol{\sigma}^*\|_2\le\|\hat{\mathbf{M}}-\mathbf{M}^*\|_F$。RII $\rho=f(\sigma_1,\sigma_2)=1-\sigma_1^2/(\sigma_1^2+\sigma_2^2)$ 有 $\|\nabla f\|_2=\frac{2\sigma_1\sigma_2}{(\sigma_1^2+\sigma_2^2)^{3/2}}\le2/\sigma_1$（利用 $\sigma_2\le\sigma_1$），故由中值定理
\begin{equation}
|\hat{\rho}-\rho^*|\le\frac{2}{\sigma_1}\|\hat{\boldsymbol{\sigma}}-\boldsymbol{\sigma}^*\|_2\le\frac{4\nu}{\sigma_1}\sqrt{\frac{C\log(4C/\delta)}{N_{\min}}}.
\end{equation}
令其 $\le\varepsilon$ 得 $N_{\min}\ge16C\nu^2\log(4C/\delta)/(\sigma_1^2\varepsilon^2)$。
\end{proof}

\subsection{定理~\ref{thm:mi_bound} 的证明（谱 MI 界）}
\label{app:mi_bound}
\begin{proof}
令 $q=P(X{=}f)$，$\bar{q}=1-q$，$P_Y=q\boldsymbol{\mu}_f+\bar{q}\boldsymbol{\mu}_r$，$\pi_{\min}=\min_y (\boldsymbol{\mu}_r)_y$。以 nat 为单位，$D_{\mathrm{KL}}\le\chi^2$，故 $I(X;Y\mid\theta')\le q\chi^2(\boldsymbol{\mu}_f\|P_Y)+\bar{q}\chi^2(\boldsymbol{\mu}_r\|P_Y)$。由于 $P_Y(y)\ge\bar{q}\pi_{\min}$，
\begin{equation}
\chi^2(\boldsymbol{\mu}_f\|P_Y)=\sum_y\frac{\bar{q}^2(\mu_f(y)-\mu_r(y))^2}{P_Y(y)}\le\frac{\bar{q}\,\|\boldsymbol{\mu}_f-\boldsymbol{\mu}_r\|_2^2}{\pi_{\min}},
\end{equation}
对称地 $\chi^2(\boldsymbol{\mu}_r\|P_Y)\le q\|\boldsymbol{\mu}_f-\boldsymbol{\mu}_r\|_2^2/\pi_{\min}$，故 $I\le\frac{2q\bar{q}}{\pi_{\min}}\|\boldsymbol{\mu}_f-\boldsymbol{\mu}_r\|_2^2$。在等范数条件 $\|\boldsymbol{\mu}_f\|_2=\|\boldsymbol{\mu}_r\|_2$ 下，$\|\boldsymbol{\mu}_f-\boldsymbol{\mu}_r\|_2^2=2\sigma_2^2$（定义~\ref{def:M}），故 $I\le\frac{4q\bar{q}\sigma_2^2}{\pi_{\min}}\le\frac{\sigma_2^2}{\pi_{\min}}=\frac{\sigma_1^2}{\pi_{\min}}\frac{\rho}{1-\rho}$（因为 $4q\bar{q}\le1$）。当 $\rho\to0$ 时等范数条件渐近成立，故无该条件时 $I=O(\rho)$。
\end{proof}

\subsection{定理~\ref{thm:finetune} 的证明（微调泄漏界）}
\label{app:finetune}

\begin{lemma}[交叉协方差界]
\label{lem:cross_cov}
假设对所有 $x$ 有 $\|\mathbf{J}_\theta z_{\theta_0}(x)\|_2 \le J_{\max}$。经过一个 $\lambda$-步后，对训练良好的分类器（$\varepsilon=O(\lambda)$），$\|\boldsymbol{\Sigma}_{fr}\|_F = O(\lambda)$。
\end{lemma}
\begin{proof}
对 $c_f\neq c_r$，$\theta_0$ 处的 logit 梯度内积为 $\langle\nabla_z\ell_f,\nabla_z\ell_r\rangle=\langle\mathbf{p}_f,\mathbf{p}_r\rangle-p_f(c_r)-p_r(c_f)=O(\varepsilon)$（分类器误差）。经过 $\lambda$-步后，$\|\Delta z\|_2\le\lambda L\|\mathbf{g}\|_2+O(\lambda^2)$，且 softmax $1$-Lipschitz 连续性给出 $|\langle\nabla_z\ell'_f,\nabla_z\ell'_r\rangle|\le|\langle\nabla_z\ell_f,\nabla_z\ell_r\rangle|+2\|\Delta z\|_2=O(\varepsilon+\lambda)$。通过链式法则 $\nabla_\theta\ell=\mathbf{J}_\theta z^\top\nabla_z\ell$ 与 $\|\mathbf{J}_\theta z\|_2\le J_{\max}$，参数空间协方差继承该尺度：对训练良好的分类器，$\|\boldsymbol{\Sigma}_{fr}\|_F=O(\varepsilon+\lambda)=O(\lambda)$。
\end{proof}

\begin{proof}[定理~\ref{thm:finetune} 的证明]
以 $\mathbf{g}=\nabla_\theta\mathcal{L}(\theta_0,D_f)$ 与 $\theta'=\theta_0+\lambda\mathbf{g}$，一阶 Taylor 展开给出
\begin{equation}
\Delta\boldsymbol{\mu}_f=\hat{\boldsymbol{\mu}}_f'-\hat{\boldsymbol{\mu}}_f=\frac{\lambda}{N_f}\sum_{x\in D_f}\mathbf{J}_{\mathrm{softmax}}(f_{\theta_0}(x))\mathbf{J}_\theta f_{\theta_0}(x)\mathbf{g}+O(\lambda^2),
\end{equation}
且对保留样本 $\|\Delta\boldsymbol{\mu}_r\|_2\le\lambda L\|\mathbf{g}\|_2$（一致 softmax Jacobian 界）。平方位移满足 $\|\hat{\boldsymbol{\mu}}_f'-\hat{\boldsymbol{\mu}}_r'\|_2^2\le\|\Delta\boldsymbol{\mu}_f\|_2^2+\|\Delta\boldsymbol{\mu}_r\|_2^2+2\|\Delta\boldsymbol{\mu}_f\|_2\|\Delta\boldsymbol{\mu}_r\|_2$（我们省略与 $\boldsymbol{\mu}_f-\boldsymbol{\mu}_r$ 的一阶交叉项：一般情形它至多为 $O(\lambda^2)$ 阶，因此不改变首阶 $O(\lambda^2)$ 界；对类级遗忘，其期望由引理~\ref{lem:cross_cov} 为 $O(\lambda^3)$，故可吸收进余项）。由 $\|\Delta\boldsymbol{\mu}_f\|_2\le\lambda L\|\mathbf{g}\|_2\sqrt{\mathrm{Tr}(\boldsymbol{\Sigma}_f)}$（在小步条件 $\lambda\|\mathbf{g}\|_2=O(1)$ 下；等价地，可将上升方向归一化并把该因子吸收进 $\lambda$），
\begin{equation}
\|\hat{\boldsymbol{\mu}}_f'-\hat{\boldsymbol{\mu}}_r'\|_2^2\le\lambda^2L^2\mathrm{Tr}(\boldsymbol{\Sigma}_f)+O(\lambda^3).
\end{equation}
利用 $\|\boldsymbol{\mu}_f-\boldsymbol{\mu}_r\|_2=\sqrt{2}\sigma_2$ 与 $\rho\approx\|\boldsymbol{\mu}_f-\boldsymbol{\mu}_r\|_2^2/(2\sigma_1^2)$（$\rho\ll1$ 时）得 $\rho(\theta')\le\frac{\lambda^2L^2}{2\sigma_1^2}\mathrm{Tr}(\boldsymbol{\Sigma}_f)+O(\lambda^3)$。对类级遗忘，交叉项 $\mathbb{E}[\langle\Delta\boldsymbol{\mu}_f,\Delta\boldsymbol{\mu}_r\rangle]=\lambda^2\mathrm{Tr}(\boldsymbol{\Sigma}_{fr})+O(\lambda^3)$ 由引理~\ref{lem:cross_cov} 为 $O(\lambda^3)$，保持 $O(\lambda^2)$ 阶。
\end{proof}

\subsection{定理~\ref{thm:unified} 的证明梗概（统一分解）}
\label{app:unified}
\begin{proof}
由链式法则，$I(X;R) \le I(X;Y) + I(X;R\mid Y)$。对覆盖项，Hilbert--Schmidt 分解 $f_{Y|X}(y|x) = f_Y(y) + \sum_{k\ge 2}\sigma_k\phi_k(y)\psi_k(x)$ 给出 $\chi^2(P_{Y|X=x}\|P_Y)=\int\frac{(\sum_{k\ge2}\sigma_k\phi_k(y)\psi_k(x))^2}{f_Y(y)}\,dy=\sum_{k\ge2}\sigma_k^2\psi_k(x)^2$；对 $X$ 取平均并利用 $\mathbb{E}_X[\psi_k(X)^2]=1$ 得 $\mathbb{E}_X[\chi^2(P_{Y|X}\|P_Y)] = \sum_{k\ge 2}\sigma_k^2 = \eta^2$。由于 $D_{\mathrm{KL}} \le \chi^2$，有 $I(X;Y) \le \eta^2$。对直接项，条件于 $Y=y$，$R = \alpha X + \beta y + N$。由引理~\ref{lem:fisher}（Fisher 信息展开），$I(X;R\mid Y) = \frac{\alpha^2}{2\sigma^2}\operatorname{Var}(X) + o(\alpha^2)$。求和即得分解。
\end{proof}

\begin{proof}[定理~\ref{thm:state_bridge} 的证明]
由 $\mathbf{M}_{\mathrm{sm}} = \mathbf{B}\mathbf{A}$ 与变分刻画 $\sigma_i(\mathbf{B}\mathbf{A})=\max_{\dim V=i}\min_{v\in V}\|\mathbf{B}\mathbf{A}v\|_2/\|v\|_2$，上界由 $\|\mathbf{A}v\|_2\le\tau_1\|v\|_2$ 得：$\sigma_i(\mathbf{B}\mathbf{A})\le\tau_1\sigma_i(\mathbf{B})$。对下界，右逆 $\mathbf{A}\mathbf{A}^+=\mathbf{I}_m$ 给出 $\mathbf{B}=\mathbf{B}\mathbf{A}\mathbf{A}^+$，故 $\sigma_i(\mathbf{B})\le\|\mathbf{A}^+\|_2\,\sigma_i(\mathbf{B}\mathbf{A})=\tau_m^{-1}\sigma_i(\mathbf{B}\mathbf{A})$，即 $\sigma_i(\mathbf{B}\mathbf{A})\ge\tau_m\sigma_i(\mathbf{B})$。因此 $\frac{\tau_m}{\tau_1}\epsilon \le \tilde{\epsilon} \le \frac{\tau_1}{\tau_m}\epsilon$，其中 $\epsilon=\sigma_2/\sigma_1$，给出 $|\rho_{\mathrm{sm}}-\rho_{\mathcal{S}}| \le 2(\tau_1/\tau_m-1)\rho_{\mathcal{S}} + O((\tau_1/\tau_m-1)^2)$。$L\delta_q/\|\boldsymbol{\mu}_f\|_2$ 项计入 $\mathbf{a}_s$ 的类内方差。
\end{proof}

\subsection{桥与 $\chi^2$ 控制证明}
\label{app:bridge_proofs}

\begin{proof}[定理~\ref{thm:state_mhpr} 的证明]
$\rho_{H,\mathcal{S}} = \min_{\mathbf{w}} \|\boldsymbol{\pi}_f - \sum_k w_k\boldsymbol{\pi}_{h_k}\|_2^2 / \|\boldsymbol{\pi}_f\|_2^2 \le \|\boldsymbol{\pi}_f - \bar{\boldsymbol{\pi}}_H\|_2^2 / \|\boldsymbol{\pi}_f\|_2^2$。由于 $1/\bar{\pi}_H(m) \ge 1$，有 $\|\boldsymbol{\pi}_f - \bar{\boldsymbol{\pi}}_H\|_2^2 \le \chi^2(\boldsymbol{\pi}_f\|\bar{\boldsymbol{\pi}}_H)$。结合 $\|\boldsymbol{\pi}_f\|_2^2 \ge \gamma$ 与 $\chi^2 \le \tau^2$，得 $\rho_{H,\mathcal{S}} \le \tau^2/\gamma$。
\end{proof}

\begin{proof}[定理~\ref{thm:softmax_deviation} 的证明]
分解 $\|\hat{\boldsymbol{\mu}}_f - \mathbf{P}_{\hat{\mathbf{H}}}\hat{\boldsymbol{\mu}}_f\|_2 \le \|(\mathbf{I} - \mathbf{P}_{\mathbf{H}^*})\hat{\boldsymbol{\mu}}_f\|_2 + \|(\mathbf{P}_{\mathbf{H}^*} - \mathbf{P}_{\hat{\mathbf{H}}})\hat{\boldsymbol{\mu}}_f\|_2$。在 $\boldsymbol{\mu}_f^* \in \mathrm{rowspan}(\mathbf{H}^*)$ 下，第一项化为 $\|\hat{\boldsymbol{\mu}}_f - \boldsymbol{\mu}_f^*\|_2$。对第二项，由子空间摄动理论~\cite{stewart1990matrix}，当 $\|\hat{\mathbf{H}}-\mathbf{H}^*\|_2 \le \sigma_{\min}(\mathbf{H}^*)/2$ 时，$\|\mathbf{P}_{\hat{\mathbf{H}}} - \mathbf{P}_{\mathbf{H}^*}\|_2 \le \frac{2\sqrt{r}}{\sigma_{\min}(\mathbf{H}^*)}\|\hat{\mathbf{H}}-\mathbf{H}^*\|_2$。两个误差项都由 Hoeffding 界为 $2\nu\sqrt{C\log(2(K+1)C/\delta)/N_{\min}}$。代入并应用 delta 方法得 $\rho_H = O(1/N_{\min})$。
\end{proof}

\subsection{引理~\ref{lem:fisher} 的证明（Fisher 信息展开）}
\label{app:fisher}
\begin{proof}
假设 $f_Z\in C^2$ 且 $f_Z$ 及其导数在无穷远处衰减得足够快以使下列分部积分成立。对二元 $X\in\{0,1\}$（$P(X{=}1)=q$），$f_R(r)=f_Z(r)-q\alpha f_Z'(r)+\frac{q\alpha^2}{2}f_Z''(r)+o(\alpha^2)$。令 $\delta f=f_R-f_Z$。$f_Z+\delta f$ 的熵摄动为
\begin{equation}
H(f_Z+\delta f)=H(f_Z)-\int\delta f\,\log f_Z-\frac12\int\frac{(\delta f)^2}{f_Z}+o(\|\delta f\|^2),
\end{equation}
其中 $\int\delta f=0$（质量守恒）。$O(\alpha)$ 项消失，因为 $\int f_Z'\log f_Z=[f_Z\log f_Z]-\int f_Z'=0$；分部积分得 $\int f_Z''\log f_Z=-\int(f_Z')^2/f_Z=-J(Z)$。二次项贡献 $-\frac{q^2\alpha^2}{2}J(Z)$。因此
\begin{equation}
H(R)-H(Z)=\frac{q(1-q)\alpha^2}{2}J(Z)+o(\alpha^2)=\frac{\alpha^2}{2}\mathrm{Var}(X)J(Z)+o(\alpha^2).
\end{equation}
由于 $H(R|X)=H(Z)$，$I(X;R)=H(R)-H(R|X)=\frac{\alpha^2}{2}\mathrm{Var}(X)J(Z)+o(\alpha^2)$；对高斯噪声 $J(N)=1/\sigma^2$。
\end{proof}

\subsection{MHPR 理论保证}
\label{app:mhpr}

设 $\boldsymbol{\mu}_f \in \mathbb{R}^C$ 为遗忘类的平均 softmax 向量，$\mathbf{H}_K \in \mathbb{R}^{K\times C}$ 为 $K$ 个留出类均值的矩阵。记 $\mathcal{S}_K = \mathrm{rowspan}(\mathbf{H}_K)$，$\mathbf{P}_{\mathcal{S}_K}$ 为正交投影。

\begin{proposition}[随 $K$ 的单调改进]
\label{prop:mhpr_monotone}
对 $K_1 \le K_2$，$\rho_H(K_2) \le \rho_H(K_1)$。
\end{proposition}
\begin{proof}
由于 $\mathcal{S}_{K_1} \subseteq \mathcal{S}_{K_2}$，正交投影满足 $\mathbf{P}_{\mathcal{S}_{K_2}} = \mathbf{P}_{\mathcal{S}_{K_1}} + \mathbf{P}_{\mathcal{S}_{K_2} \cap \mathcal{S}_{K_1}^\perp}$。勾股定理给出 $\|\boldsymbol{\mu}_f - \mathbf{P}_{\mathcal{S}_{K_2}}\boldsymbol{\mu}_f\|_2^2 = \|\boldsymbol{\mu}_f - \mathbf{P}_{\mathcal{S}_{K_1}}\boldsymbol{\mu}_f\|_2^2 - \|\mathbf{P}_{\mathcal{S}_{K_2} \cap \mathcal{S}_{K_1}^\perp}\boldsymbol{\mu}_f\|_2^2 \le \|\boldsymbol{\mu}_f - \mathbf{P}_{\mathcal{S}_{K_1}}\boldsymbol{\mu}_f\|_2^2$。
\end{proof}

\begin{proposition}[偏差--方差分解]
\label{prop:mhpr_bias}
在 $\mathcal{H}_0$（$\boldsymbol{\mu}_f^* \in \mathcal{S}_K^*$）下，$\mathbb{E}[\rho_{\mathrm{est}}] = O(C/N_{\min})$。
\end{proposition}
\begin{proof}
在 $\mathcal{H}_0$（$\boldsymbol{\mu}_f^*=\mathbf{P}_{\mathcal{S}_K^*}\boldsymbol{\mu}_f^*$）下，分子展开为 $\|(\mathbf{I}-\mathbf{P}_{\mathcal{S}_K^*})\boldsymbol{\delta}_f-\Delta_H\boldsymbol{\mu}_f^*-\Delta_H\boldsymbol{\delta}_f\|_2^2$，其中 $\boldsymbol{\delta}_f=\hat{\boldsymbol{\mu}}_f-\boldsymbol{\mu}_f^*$，$\Delta_H=\mathbf{P}_{\mathcal{S}_K}-\mathbf{P}_{\mathcal{S}_K^*}$。标准界通过次高斯集中与子空间摄动~\cite{stewart1990matrix} 给出 $\|\boldsymbol{\delta}_f\|_2=O_p(\sqrt{C/N_f})$ 与 $\|\Delta_H\|_2=O_p(\sqrt{CK/N_{\min}})$；分母以 $O_p(1/\sqrt{N_f})$ 收敛。二阶 Taylor 展开 $\rho_{\mathrm{est}}\approx\rho_H^*+\nabla\rho_H^*\,(\hat{\mathbf{z}}-\mathbf{z}^*)+\frac12(\hat{\mathbf{z}}-\mathbf{z}^*)^\top\nabla^2\rho_H^*\,(\hat{\mathbf{z}}-\mathbf{z}^*)$ 的线性项均值为零（估计量一阶无偏）；二次项为 $O(\|\hat{\mathbf{z}}-\mathbf{z}^*\|_2^2)=O_p(CK/N_{\min})$。因此 delta 方法给出 $\mathbb{E}[\rho_{\mathrm{est}}]=O(CK/N_{\min})$，当 $K=O(1)$ 时为 $O(C/N_{\min})$。
\end{proof}

\begin{proposition}[有限样本界]
\label{prop:mhpr_confidence}
对方差代理 $\nu^2$ 的次高斯 softmax 向量且 $\sigma_{\min}(\mathbf{H}^*)>0$，
\begin{equation}
|\hat{\rho}_H - \rho_H^*| \le \frac{4\nu}{\|\boldsymbol{\mu}_f^*\|_2}\left(2+\frac{\sqrt{K}\,\|\boldsymbol{\mu}_f^*\|_2}{\sigma_{\min}(\mathbf{H}^*)}\right)\sqrt{\frac{C(K+1)\log(2(K+1)C/\delta)}{N_{\min}}}
\end{equation}
置信度 $\ge 1-\delta$。
\end{proposition}
\begin{proof}
令 $\hat{\mathbf{z}}=(\hat{\boldsymbol{\mu}}_f,\hat{\boldsymbol{\mu}}_{h_1},\dots,\hat{\boldsymbol{\mu}}_{h_K})$。对遗忘方向，$\|\nabla_{\boldsymbol{\mu}_f}\phi\|_2\le4/\|\boldsymbol{\mu}_f\|_2$（由 $\nabla\phi=2(\mathbf{I}-\mathbf{P}_{\mathcal{S}_K})\boldsymbol{\mu}_f/\|\boldsymbol{\mu}_f\|_2^2-2\phi\boldsymbol{\mu}_f/\|\boldsymbol{\mu}_f\|_2^2$）；对每个留出方向，Davis--Kahan~\cite{stewart1990matrix} 给出 $\|\nabla_{\boldsymbol{\mu}_{h_k}}\phi\|_2\le2/\sigma_{\min}(\mathbf{H}^*)$。当 $\|\hat{\boldsymbol{\mu}}_f\|_2\ge\|\boldsymbol{\mu}_f^*\|_2/2$（大 $N_f$）时，组合 Lipschitz 常数为 $L=\frac{4}{\|\boldsymbol{\mu}_f^*\|_2}(2+\frac{\sqrt{K}\|\boldsymbol{\mu}_f^*\|_2}{\sigma_{\min}(\mathbf{H}^*)})$。对 $(K+1)C$ 个分量做 union bound 的 Hoeffding 给出 $\|\hat{\mathbf{z}}-\mathbf{z}^*\|_\infty\le2\nu\sqrt{\log(2(K+1)C/\delta)/N_{\min}}$（高概率），故 $|\hat{\rho}_H-\rho_H^*|\le L\sqrt{(K+1)C}\|\hat{\mathbf{z}}-\mathbf{z}^*\|_\infty$。
\end{proof}

\subsection{状态空间 MHPR 界验证}
\label{app:chi2}

我们在 MNIST 类-5 设定（$m=20$ 状态，$K=3$ 留出）上验证定理~\ref{thm:state_mhpr}。界 $\rho_{H,\mathcal{S}} \le \chi^2/\|\boldsymbol{\pi}_f\|_2^2$ 对\emph{每一个}梯度步都成立（表~\ref{tab:chi2}）；$\rho_{H,\mathcal{S}}$ 与 $\chi^2$ 都随遗忘强度单调增长，且该界保守（比值 $\sim5\times10^2$--$2\times10^3$），符合最坏情况不等式的预期。

\begin{table}[t]
\centering
\small
\caption{状态空间 MHPR 界验证（MNIST，$m{=}20$，$K{=}3$）。}\label{tab:chi2}
\setlength{\tabcolsep}{6pt}
\begin{tabular}{l c c c}
\toprule
步数 & $\rho_{H,\mathcal{S}}$ & $\chi^2$ & $\chi^2/\|\boldsymbol{\pi}_f\|_2^2$ \\
\midrule
0     & 0.999758 & 272.7 & $5.20\times 10^2$ \\
1     & 0.999758 & 273.0 & $5.20\times 10^2$ \\
3     & 0.999758 & 279.0 & $5.30\times 10^2$ \\
5     & 0.999753 & 279.1 & $5.30\times 10^2$ \\
10    & 0.999752 & 302.3 & $5.75\times 10^2$ \\
20    & 0.999746 & 346.2 & $6.50\times 10^2$ \\
50    & 0.999613 & 558.4 & $1.04\times 10^3$ \\
100   & 0.998588 & 1095.8 & $2.04\times 10^3$ \\
\bottomrule
\end{tabular}
\end{table}

\subsection{PSSD 证明}

\begin{proof}[命题~\ref{prop:pssd_rii} 的证明]
$\Delta_f = \frac1{N_f}\sum_{x\in D_f}(1-\pi_r(s(x))) = 1-\sum_s\pi_f(s)\pi_r(s)=1-\langle\boldsymbol{\pi}_f,\boldsymbol{\pi}_r\rangle$。对差距项，当 $\|\boldsymbol{\pi}_f\|_2=\|\boldsymbol{\pi}_r\|_2=s$ 时，状态空间 RII 为 $\rho_{\mathcal{S}}=\frac{s-\langle\boldsymbol{\pi}_f,\boldsymbol{\pi}_r\rangle}{2s}$。围绕 $\boldsymbol{\pi}_f=\boldsymbol{\pi}_r$ 展开得 $|\Delta_f-\rho_{\mathcal{S}}|\le\|\boldsymbol{\pi}_f-\boldsymbol{\pi}_r\|_1(1-1/m)+O(\|\boldsymbol{\pi}_f-\boldsymbol{\pi}_r\|_2^2)$，其中 $\ell_1$ 项通过 $|\langle\boldsymbol{\pi}_f,\boldsymbol{\pi}_r\rangle-\|\boldsymbol{\pi}_r\|_2^2|\le\|\boldsymbol{\pi}_f-\boldsymbol{\pi}_r\|_1\|\boldsymbol{\pi}_r\|_\infty\le\|\boldsymbol{\pi}_f-\boldsymbol{\pi}_r\|_1(1-1/m)$ 占主导。
\end{proof}

\begin{proof}[命题~\ref{prop:pssd_conc} 的证明]
每个 $\delta(x_i)\in[0,1]$ 是有界随机变量。Hoeffding 给出 $\mathbb{P}(|\hat{\Delta}_f-\Delta_f|\ge\varepsilon) \le 2\exp(-2N_f\varepsilon^2)$。对 $\hat{\Psi}=\hat{\Delta}_f-\hat{\Delta}_r$，两样本 Hoeffding 界给出 $\mathbb{P}(|\hat{\Psi}-\Psi|\ge\varepsilon) \le 4\exp(-N_{\min}\varepsilon^2/2)$。对 top-$k$，Bonferroni 给出 $\mathbb{P}(|\hat{\Delta}_{(k)}-\Delta_{(k)}|\ge\varepsilon) \le 2k\exp(-2N_f\varepsilon^2/k^2)$。在置信度 $1-\delta$ 下解出 $\varepsilon$ 即得所述速率。与命题~\ref{prop:confidence} 不同，无需次高斯假设。
\end{proof}

\begin{proof}[命题~\ref{prop:ps_mhpr} 的证明]
对固定子空间 $\mathcal{S}_H$，残差 $\|\mathbf{v}-\mathbf{P}_{\mathcal{S}_H}\mathbf{v}\|_2^2$ 是关于 $\mathbf{v}$ 的凸二次型。Jensen 给出
\[
\frac1{N_f}\sum_{x\in D_f}\rho_H(x) \ge \rho_H\!\left(\frac1{N_f}\sum_{x\in D_f}\mathbf{e}_{s(x)}\right) = \rho_H(\boldsymbol{\pi}_f) = \|\boldsymbol{\pi}_f\|_2^2\,\rho_{H,\mathcal{S}},
\]
因为 $\rho_H(\boldsymbol{\pi}_f)=\min_{\mathbf{w}}\|\boldsymbol{\pi}_f-\sum_k w_k\boldsymbol{\pi}_{h_k}\|_2^2=\|\boldsymbol{\pi}_f\|_2^2\rho_{H,\mathcal{S}}$。精确差距为 $\frac1{N_f}\sum_x\|(\mathbf{I}-\mathbf{P}_{\mathcal{S}_H})(\mathbf{e}_{s(x)}-\boldsymbol{\pi}_f)\|_2^2\ge0$，仅当遗忘类是状态齐次（$\boldsymbol{\pi}_f$ 为 Kronecker delta）时消失。
\end{proof}

\begin{proof}[定理~\ref{thm:pssd_mia} 的证明]
(i) 若 $\Psi>0$，$\delta$ 的遗忘侧均值超过保留侧均值；Bayes 最优检验是似然比检验，故存在阈值使 $\mathrm{TPR}>\mathrm{FPR}$，即 $\mathrm{AUC}>0.5$。

(ii) 由 $\mathrm{H}^2(P,Q)=1-\sum_i\sqrt{p_iq_i}$，总变分距离满足 $\mathrm{TV}(P_{\delta|f},P_{\delta|r})\le\mathrm{H}(P_{\delta|f},P_{\delta|r})$（Le Cam）。对阈值分类器，$\mathrm{AUC}=\frac12(1+\mathrm{TPR}-\mathrm{FPR})\ge\frac12(1-\mathrm{TV})$，因为在最优阈值处 $\mathrm{TPR}-\mathrm{FPR}\ge1-\mathrm{TV}$。综合即得 Hellinger 界 $\mathrm{AUC}\ge1-\frac12\mathrm{H}^2$。

(iii) AUC 是核有界的 $U$-统计量，故 Hoeffding 给出 $|\widehat{\mathrm{AUC}}-\mathrm{AUC}|\le\sqrt{\log(2/\zeta)/(2(N_f+N_r))}$。

对基于 softmax 的 MIA（表~\ref{tab:pssd_results}）的经验优势源于 $\delta(x)$ 通过 $\pi_r(Q(z(x)))$ 依赖 logit 几何，而后者在 softmax 概率弥散时仍可判别。
\end{proof}

\subsection{引理~\ref{lem:quantization} 的证明（量化误差界）}
\begin{proof}
对独立同分布、$\|z\|_2\le R$ 且最优 $m$ 中心失真 $D^*=\mathbb{E}[\min_{c\in\mathcal{C}^*}\|z-c\|_2^2]$ 的向量，Kumar--Kannan~\cite{kumar2010clustering} 的过量失真界 $D(\hat{\mathcal{C}})-D^*=O(R^2\sqrt{Cm/N})$ 高概率成立。由于 $\mathbb{E}[\|z-\tilde{z}\|_2]\le\sqrt{D(\hat{\mathcal{C}})}$ 且 $D^*\le R^2$，开方得 $\tilde{O}(R\sqrt{Cm/N})$；$O(R\sqrt{\log(1/\zeta)/N})$ 尾界来自对经验失真估计的 union bound。
\end{proof}

\subsection{温度缩放信息保持}
\label{app:temp}
\begin{proposition}[保序]
\label{prop:rank_preserve}
设 $\theta_1',\theta_2'$ 为参数状态，$\rho_T(\theta')$ 为 $\mathbf{p}_T(x)=\mathrm{Softmax}(f(x)/T)$ 在温度 $T$ 下的 RII。若 $\rho_1(\theta_1') \le \rho_1(\theta_2')$，则存在 $T_0\ge1$ 使对所有 $T\ge T_0$ 有 $\rho_T(\theta_1') \le \rho_T(\theta_2')$，前提是保留类均值保持良态。
\end{proposition}
\begin{proof}
当 $T\to\infty$ 时，$\mathbf{p}_T\to\mathbf{1}/C$ 且 $\rho_T\to0$；收敛速率由 $\|\boldsymbol{\mu}_f^{(T)}-\boldsymbol{\mu}_r^{(T)}\|_2\propto T^{-1}\|\boldsymbol{\mu}_f^{(1)}-\boldsymbol{\mu}_r^{(1)}\|_2$ 控制。由于 $T=1$ 的排序是严格的，乘法缩放对一切充分大的 $T$ 保持该排序。
\end{proof}

% =============================================================================
% 附录 D 实验细节与 MIA 协议
% =============================================================================
\section{实验细节与 MIA 协议}
\label{app:details}

\subsection{视觉基准}

表~\ref{tab:exp_vision} 汇总视觉基准配置；类级基准基于三个种子（0、1、2）运行并以均值 $\pm$ 标准差报告，无数据增强、逐通道归一化（CIFAR-10：均值 $(0.4914,0.4822,0.4465)$，标准差 $(0.2023,0.1994,0.2010)$；Fashion-MNIST/MNIST：均值 $0.2860$，标准差 $0.3530$）。

\begin{table}[t]
\centering
\small
\caption{视觉基准配置。``类级''指第~\ref{sec:benchmark}~节的协议。}\label{tab:exp_vision}
\setlength{\tabcolsep}{2pt}
\begin{tabular}{l l c c c c}
\toprule
数据集 & 架构 & 优化器 & lr & batch & epoch \\
\midrule
MNIST & SimpleMLP（784$\to$128$\to$10） & Adam & $10^{-3}$ & 64 & 10 \\
Fashion-MNIST & MLP & Adam & $10^{-3}$ & 64 & 10 \\
CIFAR-10 & SmallCNN（3 卷积 + 2 全连接） & Adam & $10^{-3}$ & 64 & 10 \\
CIFAR-100 & SmallCNN / ResNet-18 & Adam & $10^{-3}$ & 64 & 10 \\
MNIST（过拟合） & DeepMLP & Adam & $10^{-3}$ & 64 & 50 \\
\bottomrule
\end{tabular}
\end{table}

随机遗忘基线（表~\ref{tab:main}）丢弃随机 5\% 子集；类级协议在 $\{0,\dots,6\}$ 上训练，遗忘类~3、留出类 $\{7,8,9\}$（CIFAR-10 与 Fashion-MNIST），并在 20 类 CIFAR-100 子集上训练，遗忘类~3、留出 $\{20,21,22\}$。所有实验在 Apple Silicon（M5 Pro，24~GB 统一内存）上以 MPS 加速运行；最大的视觉模型（CIFAR-100 上的 ResNet-18）几分钟内即可训练完成，无需 GPU 集群。

\subsection{语言模型实验}

表~\ref{tab:exp_llm} 汇总 LLM 配置。TOFU~\cite{maini2024tofu} 作者级协议使用一个遗忘作者（40 个 QA 对）与 90 作者保留切分；评估使用 20 个遗忘与 40 个保留 QA 对，最大序列长度 384，并在三个随机评估子采样上重复。LoRA（rank 8，scale 20，dropout 0）应用于 32 层中的 8 层。

\begin{table}[t]
\centering
\small
\caption{语言模型配置。}\label{tab:exp_llm}
\setlength{\tabcolsep}{2pt}
\begin{tabular}{l l l l}
\toprule
实验 & 模型 & PEFT & 遗忘配方 \\
\midrule
TOFU 7B & LLaMA-2-7B（TOFU 检查点） & LoRA $r{=}8$，scale 20，8/32 层 & NoUnlearn：原检查点；FineTune：AdamW $10^{-5}$，60 步，batch 1；NegGrad：SGD $2{\times}10^{-5}$，10 步；Retrain：基座 + LoRA on retain，60 步 \\
MUSE 7B & LLaMA-2-7B（基座） & LoRA $r{=}8$，scale 20，8/32 层 & Teacher：遗忘书籍上 200 步 LoRA；FineTune：AdamW $10^{-5}$，60 步；NegGrad：SGD $2{\times}10^{-5}$，15 步；Retrain：基座 + LoRA on retain，60 步 \\
TOFU 0.5B & Qwen2.5-0.5B & 全量微调 & AdamW $5{\times}10^{-5}$；FineTune 40 步，NegGrad 8 步，batch 4 \\
AG News & DistilBERT & 全量微调 & Adam $2{\times}10^{-5}$，2 epoch，batch 32；每类 4000 训练 / 1500 评估；遗忘类 = World 新闻 \\
\bottomrule
\end{tabular}
\end{table}

NegGrad 强度扫描设置如第~\ref{sec:llm}~节所述。

\subsection{成员推理攻击细节}
\label{app:mia}

本文出现两种不同的 MIA 协议，如各表标题与第~\ref{sec:experiments}~节设置中的``关于 MIA 的记号''段落所述。

\textbf{损失阈值 AUC（类级基准，表~\ref{tab:benchmark}）。}对每个遗忘与保留样本计算模型交叉熵损失；每个样本的得分为其损失，报告的 MIA-loss 是对这些得分取阈值得到的 AUC（遗忘样本为正类）。不涉及影子模型或重训练：直接评估目标模型。由于遗忘类在训练协议中是未见类，其损失超过保留样本，产生第~\ref{sec:benchmark}~节讨论的``反向''排序（oracle 达到 AUC${\to}0$，这本身就是成功擦除的证据）。

\textbf{训练/测试成员准确率（随机基线，表~\ref{tab:main}）。}我们按 softmax 置信度（或等价地，负损失）给每个样本打分，并通过阈值分类训练/测试成员性；报告值为两类的平均成员准确率。这度量一般记忆（训练/测试置信度差距），与遗忘/保留不可区分性正交（第~\ref{sec:mia}~节）。

\textbf{为何简单攻击足够。}本文目标是排序遗忘方法而非最大化攻击成功率，且证书 RII 是分布量而非对抗界。更强攻击（如 LiRA 或 Carlini 等人的一阶原理分析~\cite{carlini2022membership}）需要影子模型或重复训练；在统一单模型协议下，额外威力不会改变方法的相对排序，因为任何黑盒攻击者可用的输出级信号就是简单阈值攻击所用的逐样本损失或置信度。因此我们报告简单攻击，并把更强攻击的比较留待未来工作。

% =============================================================================
% 参考文献（保留英文原文）
% =============================================================================
\begin{thebibliography}{99}
\bibitem{bourtoule2021machine}
Bourtoule L, Chandrasekaran V, Choquette-Choo C A, et al (2021) Machine unlearning. In: 2021 IEEE Symposium on Security and Privacy (SP), IEEE

\bibitem{nguyen2025survey}
Nguyen T T, Huynh T T, Ren Z, et al (2025) A survey of machine unlearning. ACM Transactions on Intelligent Systems and Technology 16(5):1-46

\bibitem{cosma2026ruler}
Cosma G, Finke A (2026) RULER: representation-level verification of machine unlearning. arXiv:2605.27569

\bibitem{yong2026erased}
Yong T P Y, Ong W K, Chan C S (2026) Erased, but not gone: output forgetting is not true forgetting. arXiv:2606.25001

\bibitem{sanh2020distilbert}
Sanh V, Debut L, Chaumond J, Wolf T (2019) DistilBERT, a distilled version of BERT: smaller, faster, cheaper and lighter. arXiv:1910.01108

\bibitem{qwen2024qwen2}
Qwen Team (2024) Qwen2.5 technical report. arXiv:2412.15115

\bibitem{maini2024tofu}
Maini P, Feng Z, Schwarzschild A, Lipton Z C, Kolter J Z (2024) TOFU: a task of fictitious unlearning for LLMs. In: Proceedings of the First Conference on Language Modeling (COLM)

\bibitem{shokri2017membership}
Shokri R, Stronati M, Song C, Shmatikov V (2017) Membership inference attacks against machine learning models. In: 2017 IEEE Symposium on Security and Privacy (SP), IEEE

\bibitem{guo2020certified}
Guo C, Goldstein T, Hannun A, van der Maaten L (2020) Certified data removal from machine learning models. In: Proceedings of the 37th International Conference on Machine Learning (ICML), PMLR

\bibitem{sommer2022towards}
Sommer D M, Song L, Wagh S, Mittal P (2022) Athena: probabilistic verification of machine unlearning. Proceedings on Privacy Enhancing Technologies 2022(3):268-290

\bibitem{ye2026auditing}
Ye D, Zhu T, Huang R, et al (2026) Auditing machine unlearning: a systematic research on whether models truly forget. arXiv:2606.16110

\bibitem{cao2015towards}
Cao Y, Yang J (2015) Towards making systems forget with machine unlearning. In: 2015 IEEE Symposium on Security and Privacy (SP), IEEE

\bibitem{golatkar2020eternal}
Golatkar A, Achille A, Soatto S (2020) Eternal sunshine of the spotless net: selective forgetting in deep networks. In: 2020 IEEE/CVF Conference on Computer Vision and Pattern Recognition (CVPR), IEEE

\bibitem{ginart2019making}
Ginart A, Guan M Y, Valiant G, Zou J (2019) Making AI forget you: data deletion in machine learning. In: Advances in Neural Information Processing Systems 32 (NeurIPS)

\bibitem{sekhari2021remember}
Sekhari A, Acharya J, Kamath G, Suresh A T (2021) Remember what you want to forget: algorithms for machine unlearning. In: Advances in Neural Information Processing Systems 34 (NeurIPS)

\bibitem{carlini2022membership}
Carlini N, Chien S, Nasr M, et al (2022) Membership inference attacks from first principles. In: 2022 IEEE Symposium on Security and Privacy (SP), IEEE

\bibitem{zhang2024fragile}
Zhang B, Chen Z, Shen C, Li J (2024) Verification of machine unlearning is fragile. In: Proceedings of the 41st International Conference on Machine Learning (ICML), PMLR

\bibitem{zheng2026bias}
Zheng W, Chen K, Guo Y, Xiao Y (2026) Classification-head bias in class-level machine unlearning: diagnosis, mitigation, and evaluation. arXiv:2605.08730

\bibitem{xue2026verification}
Xue L, Hu S, Lu W, et al (2026) Towards reliable forgetting: a survey on machine unlearning verification. ACM Computing Surveys 58(12), Article 314. doi:10.1145/3807451

\bibitem{shannon1949communication}
Shannon C E (1949) Communication theory of secrecy systems. Bell System Technical Journal 28(4):656-715

\bibitem{bloch2021overview}
Bloch M, G\"unl\"u O, Yener A, et al (2021) An overview of information-theoretic security and privacy: metrics, limits and applications. IEEE Journal on Selected Areas in Information Theory 2(1)

\bibitem{shi2024muse}
Shi W, Lee J, Huang Y, et al (2024) MUSE: machine unlearning six-way evaluation for language models. arXiv:2407.06460

\bibitem{vershynin2018high}
Vershynin R (2018) High-Dimensional Probability: An Introduction with Applications in Data Science. Cambridge University Press, Cambridge

\bibitem{stewart1990matrix}
Stewart G W, Sun J (1990) Matrix Perturbation Theory. Academic Press, Boston

\bibitem{kumar2010clustering}
Kumar A, Kannan R (2010) Clustering with spectral norm and the k-means algorithm. In: 2010 IEEE 51st Annual Symposium on Foundations of Computer Science (FOCS), IEEE
\end{thebibliography}


\end{document}
